% Options for packages loaded elsewhere
\PassOptionsToPackage{unicode}{hyperref}
\PassOptionsToPackage{hyphens}{url}
\PassOptionsToPackage{dvipsnames,svgnames,x11names}{xcolor}
\documentclass[
  10pt,
]{article}
\usepackage{xcolor}
\usepackage[margin=2.2cm]{geometry}
\usepackage{amsmath,amssymb}
\setcounter{secnumdepth}{-\maxdimen} % remove section numbering
\usepackage{iftex}
\ifPDFTeX
  \usepackage[T1]{fontenc}
  \usepackage[utf8]{inputenc}
  \usepackage{textcomp} % provide euro and other symbols
\else % if luatex or xetex
  \usepackage{unicode-math} % this also loads fontspec
  \defaultfontfeatures{Scale=MatchLowercase}
  \defaultfontfeatures[\rmfamily]{Ligatures=TeX,Scale=1}
\fi
\usepackage{lmodern}
\ifPDFTeX\else
  % xetex/luatex font selection
\fi
% Use upquote if available, for straight quotes in verbatim environments
\IfFileExists{upquote.sty}{\usepackage{upquote}}{}
\IfFileExists{microtype.sty}{% use microtype if available
  \usepackage[]{microtype}
  \UseMicrotypeSet[protrusion]{basicmath} % disable protrusion for tt fonts
}{}
\makeatletter
\@ifundefined{KOMAClassName}{% if non-KOMA class
  \IfFileExists{parskip.sty}{%
    \usepackage{parskip}
  }{% else
    \setlength{\parindent}{0pt}
    \setlength{\parskip}{6pt plus 2pt minus 1pt}}
}{% if KOMA class
  \KOMAoptions{parskip=half}}
\makeatother
\usepackage{color}
\usepackage{fancyvrb}
\newcommand{\VerbBar}{|}
\newcommand{\VERB}{\Verb[commandchars=\\\{\}]}
\DefineVerbatimEnvironment{Highlighting}{Verbatim}{commandchars=\\\{\}}
% Add ',fontsize=\small' for more characters per line
\newenvironment{Shaded}{}{}
\newcommand{\AlertTok}[1]{\textcolor[rgb]{1.00,0.00,0.00}{\textbf{#1}}}
\newcommand{\AnnotationTok}[1]{\textcolor[rgb]{0.38,0.63,0.69}{\textbf{\textit{#1}}}}
\newcommand{\AttributeTok}[1]{\textcolor[rgb]{0.49,0.56,0.16}{#1}}
\newcommand{\BaseNTok}[1]{\textcolor[rgb]{0.25,0.63,0.44}{#1}}
\newcommand{\BuiltInTok}[1]{\textcolor[rgb]{0.00,0.50,0.00}{#1}}
\newcommand{\CharTok}[1]{\textcolor[rgb]{0.25,0.44,0.63}{#1}}
\newcommand{\CommentTok}[1]{\textcolor[rgb]{0.38,0.63,0.69}{\textit{#1}}}
\newcommand{\CommentVarTok}[1]{\textcolor[rgb]{0.38,0.63,0.69}{\textbf{\textit{#1}}}}
\newcommand{\ConstantTok}[1]{\textcolor[rgb]{0.53,0.00,0.00}{#1}}
\newcommand{\ControlFlowTok}[1]{\textcolor[rgb]{0.00,0.44,0.13}{\textbf{#1}}}
\newcommand{\DataTypeTok}[1]{\textcolor[rgb]{0.56,0.13,0.00}{#1}}
\newcommand{\DecValTok}[1]{\textcolor[rgb]{0.25,0.63,0.44}{#1}}
\newcommand{\DocumentationTok}[1]{\textcolor[rgb]{0.73,0.13,0.13}{\textit{#1}}}
\newcommand{\ErrorTok}[1]{\textcolor[rgb]{1.00,0.00,0.00}{\textbf{#1}}}
\newcommand{\ExtensionTok}[1]{#1}
\newcommand{\FloatTok}[1]{\textcolor[rgb]{0.25,0.63,0.44}{#1}}
\newcommand{\FunctionTok}[1]{\textcolor[rgb]{0.02,0.16,0.49}{#1}}
\newcommand{\ImportTok}[1]{\textcolor[rgb]{0.00,0.50,0.00}{\textbf{#1}}}
\newcommand{\InformationTok}[1]{\textcolor[rgb]{0.38,0.63,0.69}{\textbf{\textit{#1}}}}
\newcommand{\KeywordTok}[1]{\textcolor[rgb]{0.00,0.44,0.13}{\textbf{#1}}}
\newcommand{\NormalTok}[1]{#1}
\newcommand{\OperatorTok}[1]{\textcolor[rgb]{0.40,0.40,0.40}{#1}}
\newcommand{\OtherTok}[1]{\textcolor[rgb]{0.00,0.44,0.13}{#1}}
\newcommand{\PreprocessorTok}[1]{\textcolor[rgb]{0.74,0.48,0.00}{#1}}
\newcommand{\RegionMarkerTok}[1]{#1}
\newcommand{\SpecialCharTok}[1]{\textcolor[rgb]{0.25,0.44,0.63}{#1}}
\newcommand{\SpecialStringTok}[1]{\textcolor[rgb]{0.73,0.40,0.53}{#1}}
\newcommand{\StringTok}[1]{\textcolor[rgb]{0.25,0.44,0.63}{#1}}
\newcommand{\VariableTok}[1]{\textcolor[rgb]{0.10,0.09,0.49}{#1}}
\newcommand{\VerbatimStringTok}[1]{\textcolor[rgb]{0.25,0.44,0.63}{#1}}
\newcommand{\WarningTok}[1]{\textcolor[rgb]{0.38,0.63,0.69}{\textbf{\textit{#1}}}}
\usepackage{longtable,booktabs,array}
\usepackage{caption}
\captionsetup[table]{skip=6pt}
\newcounter{none} % for unnumbered tables
\usepackage{calc} % for calculating minipage widths
% Correct order of tables after \paragraph or \subparagraph
\usepackage{etoolbox}
\makeatletter
\patchcmd\longtable{\par}{\if@noskipsec\mbox{}\fi\par}{}{}
\makeatother
\usepackage{graphicx}
\makeatletter
\newsavebox\pandoc@box
\newcommand*\pandocbounded[1]{% scales image to fit in text height/width
  \sbox\pandoc@box{#1}%
  \Gscale@div\@tempa{\textheight}{\dimexpr\ht\pandoc@box+\dp\pandoc@box\relax}%
  \Gscale@div\@tempb{\linewidth}{\wd\pandoc@box}%
  \ifdim\@tempb\p@<\@tempa\p@\let\@tempa\@tempb\fi% select the smaller of both
  \ifdim\@tempa\p@<\p@\scalebox{\@tempa}{\usebox\pandoc@box}%
  \else\usebox{\pandoc@box}%
  \fi%
}
% Set default figure placement to htbp
\def\fps@figure{htbp}
\makeatother
\setlength{\emergencystretch}{3em} % prevent overfull lines
\providecommand{\tightlist}{%
  \setlength{\itemsep}{0pt}\setlength{\parskip}{0pt}}
\usepackage{fontspec}
\setmainfont{STIX Two Text}
\setmonofont{Menlo}[Scale=0.85]
\usepackage{xeCJK}
\setCJKmainfont{Songti SC}
\usepackage{microtype}
\usepackage{etoolbox}
\AtBeginEnvironment{longtable}{\footnotesize}
\usepackage{xurl}
\setlength{\emergencystretch}{3em}
\sloppy
\setCJKmonofont{Songti SC}
\usepackage{unicode-math}
\setmathfont{STIX Two Math}
\usepackage{newunicodechar}
\newunicodechar{→}{\ensuremath{\rightarrow}}
\newunicodechar{∎}{\ensuremath{\blacksquare}}
\newunicodechar{≠}{\ensuremath{\neq}}
\newunicodechar{≈}{\ensuremath{\approx}}
\newunicodechar{↔}{\ensuremath{\leftrightarrow}}
\newunicodechar{≪}{\ensuremath{\ll}}
\newunicodechar{↛}{\ensuremath{\nrightarrow}}
\newunicodechar{・}{\textperiodcentered}
\newunicodechar{⇒}{\ensuremath{\Rightarrow}}
\newunicodechar{∪}{\ensuremath{\cup}}
\newunicodechar{≥}{\ensuremath{\geq}}
\newunicodechar{≤}{\ensuremath{\leq}}
\newunicodechar{⟨}{\ensuremath{\langle}}
\newunicodechar{⟩}{\ensuremath{\rangle}}
\usepackage{bookmark}
\IfFileExists{xurl.sty}{\usepackage{xurl}}{} % add URL line breaks if available
\urlstyle{same}
% fallback for those not using the hyperref driver hyperxmp:
\makeatletter
\@ifundefined{xmpquote}{\newcommand{\xmpquote}[1]{#1}}{}
\makeatother
\hypersetup{
  pdftitle={Metered Commit Gates: Energy-to-Correct-Completion for Agentic Systems},
  pdfauthor={David Charlot, Open Interface Engineering},
  colorlinks=true,
  linkcolor={blue},
  filecolor={Maroon},
  citecolor={Blue},
  urlcolor={blue},
  pdfcreator={LaTeX via pandoc}}

\title{Metered Commit Gates: Energy-to-Correct-Completion for Agentic
Systems}
\author{David Charlot, Open Interface Engineering}
\date{Version 1.1, 11 October 2026. HTML:
https://research.openie.dev/papers/gates/}

\begin{document}
\maketitle

{
\setcounter{tocdepth}{2}
\tableofcontents
}
\emph{Version 1.2, 11 October 2026. Data and code:
\href{https://research.openie.dev/data/\#gates}{research.openie.dev/data/}.
LaTeX PDF:
\href{https://research.openie.dev/pdfs/latex/gates.pdf}{/pdfs/latex/gates.pdf}
(\href{https://research.openie.dev/pdfs/latex/gates.tex}{source}).}

\subsection{0. Overview}\label{overview}

This paper defines a gate that lets a machine act only after a check
passes, and a rule that tries the cheapest method first. It reports
software results and one board-level energy measurement, each labeled by
how it was obtained.

\textbf{In plain words}

\begin{enumerate}
\def\labelenumi{\arabic{enumi}.}
\tightlist
\item
  A computer should not take an action that cannot be undone until a
  simple check says it is safe.
\item
  It should try the cheapest way to finish a job first and use a big AI
  model only if nothing cheaper works.
\item
  Every decision is written down with the energy it used, so anyone can
  check it later.
\end{enumerate}

This paper has two parts that share one rule: correctness is the gate,
and joules are the ranking. Part I is the commit gate. A proposal may
not authorize itself. An irreversible act runs only behind a
certificate, and every decision, commit or refuse, is written to a
receipt. Part II is Mixture of Limits: a task needs a minimal expected
energy to close, and energy spent after the task is complete buys no
verified progress. Its standard instance is a cascade that tries the
cheapest gear first: Lookup, then Formula, then Solver, then Model. The
cascade stops as soon as the task completion predicate \(C(z)\) holds,
and it refuses with a receipt when the next gear would break the budget.
Among systems that complete the task correctly, the one that needs fewer
joules ranks higher.

One board measurement anchors the gate fabric. On an AMD Kria KV260, the
onboard INA260 read SOM-level A/B/A differences of 11,428 pJ per gate
cycle (standard error 129 pJ, N = 4) for the WCA gate fabric, and 98.0
pJ per seed cycle (standard error 30.8 pJ, N = 4) for the seed fabric.
The INA260 sits on \texttt{VCC\_SOM}, which feeds the processing system,
the programmable logic, DDR memory and the module regulators
(\href{https://gbmc.googlesource.com/linux/+/ccd1cdca5cd433c8a5dff78b69a79b31d9b77ee1/arch/arm64/boot/dts/xilinx/zynqmp-sck-kv-g-revB.dtso}{KV260
device tree, \texttt{ina260@40}}), so these are whole-module
differences, not a core-rail or PL-only figure (Section I.3.7). The
seed-fabric difference is about 1 mW at 5 V, near the part's 1.25 mA
current resolution, and is reported as inconclusive until a PL-rail
meter and at least 10 passes are run.

\subsubsection{0.1 Contributions}\label{contributions}

\begin{enumerate}
\def\labelenumi{\arabic{enumi}.}
\tightlist
\item
  A commit envelope: proposal, then certificate, then commit or refuse,
  then receipt. The certificate is conjunctive: look-up allow AND an
  energy predicate, optionally AND a control barrier function. Refusal
  reasons are typed. A Rust software reference implementation, Wise
  Computer Automation (WCA), runs it (Part I).
\item
  A typed provenance invariant. Output labeled \texttt{ModelGenerated}
  can never be promoted to \texttt{Deterministic} (Parts I and II).
\item
  A cascade with named refusals, and the condition under which cheapest
  first is optimal in expected joules, which follows from Simon and
  Kadane's ordering result for sequential search: try gears in
  decreasing order of \(p_i / c_i\), where \(p_i\) is the probability
  that gear \(i\) closes the task and \(c_i\) its expected energy in
  joules (Simon and Kadane, 1975) (Part II).
\item
  Energy to correct completion as the ranking, defined in Section 0.2.
\item
  Receipts as an instrument. Each receipt holds the prediction hash and
  its timestamp before the act, the outcome, the match, the joules by
  tier (measured, reported, estimated, never summed) and the refusal
  reason.
\end{enumerate}

\subsubsection{0.2 Energy to correct
completion}\label{energy-to-correct-completion}

Let \(C(z)\) be the task completion predicate, with any deadline written
into it. For a system \(S\) with policies \(\pi\) and an allowed failure
rate \(\delta\),

\[J^*(S) = \min_{\pi} \; \mathbb{E}[J_\pi] \quad \text{subject to} \quad \Pr[C(z) = 1] \ge 1 - \delta .\]

Correctness and deadlines are constraints. Among systems that pass the
gate, joules alone order them. Ordering systems by the real number
\(J^*\) is a total preorder, since any two reals compare. If each
reading carries meter error at most \(\epsilon\) joules, two systems
whose \(J^*\) differ by more than \(2\epsilon\) keep their order. MLPerf
Power gates every power result behind an accuracy target. \(J^*\) has
the same form.

\subsubsection{0.3 How to read this paper}\label{how-to-read-this-paper}

Sections in Part I are numbered I.n and sections in Part II are numbered
II.n. Each part has its own reference list. Two figures carry the main
ideas: Figure 1 (Section I.1) draws the commit gate, and Figure 2
(Section II.1) draws the floor and the plateau of Mixture of Limits.
Prices and demand are not argued here.

\subsection{Part I. The commit gate}\label{part-i.-the-commit-gate}

\subsection{I. Abstract}\label{i.-abstract}

Agents propose actions. Irreversible work begins only when an action is
allowed to run. This paper treats that permission step as a formal
object. Linus Lee's notational intelligence is the observation that a
change in symbols can make some thoughts cheap and others expressible.
The claim here is narrower. For machines that can move matter or call an
irreversible tool, the useful notation is a runtime law: a proposal may
not authorize itself. The software reference, Wise Computer Automation
(WCA), records the law as four objects under schema
\texttt{wca.commit.v1}: a proposal, a certificate, a typed refuse
reason, and a commit decision. The certificate used in the reference is
conjunctive. A look-up table (LUT) allow bit must hold, and an energy
predicate must hold. An optional control barrier function (CBF) may be
added. Read-only calls may bypass the gate. Irreversible calls may not.

Energy numbers in this paper are products of operation counts and
analytical energy constants (an OpCounter model). They are not board
power. Field-programmable gate array (FPGA) behavior has three classes
here. Icarus Verilog simulation checks the emitted register-transfer
logic. An Alchitry Pt V2 runs the closed-loop gate and agrees with the
software on every decision. It has no meter, so its energy is unmetered.
An AMD Kria KV260 runs the gate fabric under its onboard INA260 power
monitor. Halting and running the same bitstream, interleaved A/B/A,
gives a board-level metered difference of 11,428 pJ per gate cycle
(standard error 129 pJ, N = 4) for eight parallel evaluations of LUT
allow, energy predicate and barrier at 25 MHz. The smaller seed fabric
gives 98.0 pJ per seed cycle (standard error 30.8 pJ, N = 4). These are
\texttt{measured\_j} at board level. They are not core-rail readings and
not a closed-loop plant joule cost (Section I.3.7). Stage B is available
as a browser instrument at https://research.openie.dev/living/fpga-sim/.
It compiles the same Rust gate to WebAssembly (WASM) and shows LUT state
and commit traces through the WebGPU API. That instrument is an
emulation. It is not a device-under-test (DUT) measurement and it is not
Alchitry board power.

The measured software results are bounded. On a 16-step pendulum episode
with seed 1, the reference commits 1 step and refuses 15. The analytical
episode energy is 8.6795e-10 joule. On the reported Safe fixed-point
port-Hamiltonian check, false allows against the continuous energy
oracle \texttt{Vdot\ \textless{}=\ epsilon} are 0 on the stated grids,
and the price is 79 false refuses in 209,261 decisions (Section I.4.3).
That is the safety and liveness pair on that test set (Section I.4.10).
A Model Context Protocol (MCP) demo refuses an irreversible tool call
without invoking the executor. A toy comparison with an energy-set
barrier and a discrete shield shows disagreement across safety
definitions. None of these results is a claim of silicon energy
leadership, of industrial control performance, or of a core-rail
measurement.

\subsection{I. Notation}\label{i.-notation}

Terms are defined at first scientific use below and collected here so
later sections can use the short form.

{\def\LTcaptype{none} % do not increment counter
\begin{longtable}[]{@{}
  >{\raggedright\arraybackslash}p{(\linewidth - 2\tabcolsep) * \real{0.1875}}
  >{\raggedright\arraybackslash}p{(\linewidth - 2\tabcolsep) * \real{0.8125}}@{}}
\toprule\noalign{}
\begin{minipage}[b]{\linewidth}\raggedright
Term
\end{minipage} & \begin{minipage}[b]{\linewidth}\raggedright
Definition in this paper
\end{minipage} \\
\midrule\noalign{}
\endhead
\bottomrule\noalign{}
\endlastfoot
Notational intelligence & Lee's claim that better notations can raise
what a person or system can reliably do, beyond adding tools alone. \\
Commit gate & A runtime rule that either applies a proposed action or
refuses it. The proposer does not get a second, informal channel. \\
Wise Computer Automation (WCA) & The software reference that implements
the commit gate at the boundary between a proposal and an irreversible
effect. \\
Energy-First Architecture (EFA) & The design choice to put an energy
predicate inside the permission record, not only in a later report. \\
Proposal & A candidate action \texttt{u}, state, and optional analytical
energy hints. It cannot commit. \\
Certificate & The gate record: LUT allow, energy predicate, optional CBF
predicate, reasons, and metrics. \\
Refuse reason & A typed code such as \texttt{lut\_veto},
\texttt{energy\_veto}, \texttt{cbf\_veto}, \texttt{budget\_exceeded}, or
\texttt{policy}. \\
Commit decision & The envelope that sets plant action to \texttt{u} or
to zero. \\
Look-up table (LUT) & A discrete map from a compact code to allow or
refuse. \\
Ternary look-up matrix multiplier (TLMM) & The reference's cheap ternary
table-lookup proposal path. \\
Control barrier function (CBF) & A certificate of set invariance when
its inequality holds (Ames et al., 2019). \\
Model Context Protocol (MCP) & A tool-call transport. In this paper it
is an adapter under test, not a plant certificate. \\
System One & Typed decision procedures (Laya / Jev class) that choose
among known options in software. They are proposers. \\
Port-Hamiltonian & A structured energy model used here for a
passivity-style check on toy plants. \\
Runtime verification (RV) & Checking a property on an executing trace.
Related, not identical, to the commit record. \\
Device under test (DUT) & Hardware measured under a stated workload with
a stated meter. The KV260 runs of Section I.3.7 are the only DUT
measurements in this paper. \\
Field-programmable gate array (FPGA) & Reconfigurable logic. This paper
reports Verilog simulation, a programmed Alchitry Pt V2 with unmetered
energy, and board-level KV260 meter readings (Section I.3.7). \\
WebAssembly (WASM) & A portable compilation target for the browser. The
Stage B WASM FPGA emulator is shipped at
https://research.openie.dev/living/fpga-sim/. \\
WebGPU & The browser GPU API named as the display and grid-evaluation
path for that emulator. \\
Alchitry Pt V2 & Vendor documentation identifies a Xilinx Artix-7
XC7A100T. FPGA programmed (SRAM-loaded) with the closed-loop commit
gate; energy unmetered (Section I.3.7). \\
DiffLogic & Differentiable logic gate networks trained toward Boolean
structure (Petersen et al., 2022). \\
Analytical energy estimate & \texttt{J\ =\ sum\_op\ count\_op\ *\ E\_op}
with named constants. Not board power. \\
Modeled & A quantity obtained from a stated formula or fitted
structure. \\
Simulated & A quantity obtained by executing the software or Verilog
model. \\
Board-measured & A quantity from a meter on a stated device and
workload. Only the KV260 readings of NI-9 are in this class. \\
Oracle-safe & Agreement with the continuous energy oracle
\texttt{Vdot\ \textless{}=\ epsilon} on the stated plant. \\
Check-safe & The discrete Safe fixed-point / Q16.16 residual that
approximates that oracle in the gate. \\
Set-safe & Set-invariance under a barrier or shield predicate (for
example an energy-set CBF), which can disagree with oracle-safe. \\
\end{longtable}
}

\subsubsection{I. Terms used here}\label{i.-terms-used-here}

This paper already distinguishes three safety predicates.
\textbf{Oracle-safe} means the continuous
\texttt{Vdot\ \textless{}=\ epsilon} oracle. \textbf{Check-safe} means
the discrete Safe Q16.16 residual that tracks that oracle in the gate.
\textbf{Set-safe} means set-invariance under a barrier or shield. NI-4
shows check-safe and set-safe can disagree. The software reference is
the Rust implementation that produces every simulated result here.
Evidence classes are defined in the Notation table: modeled, simulated,
board-measured. The page flag \texttt{board\_metered} is true because of
the KV260 readings alone.

\textbf{Analytical constants.} \texttt{E\_LUT\_READ\ =\ 1e-12} J and the
other OpCounter constants in Section I.3.4 are engineering stand-ins for
relative comparison inside this model. They are not calibrated board
measurements. Do not read them as silicon joules. The zero-false-allows
counts against the named oracle stay as reported.

\textbf{Living figures.} Commit stack:
\href{https://research.openie.dev/living/ni/\#ni-diag-02}{/living/ni/\#ni-diag-02}.
Seed-1 runner:
\href{https://research.openie.dev/living/ni/\#ni-run-01}{/living/ni/\#ni-run-01}.
Safe false-allow=0:
\href{https://research.openie.dev/living/ni/\#ni-run-02}{/living/ni/\#ni-run-02}.
Stage B simulator:
\href{https://research.openie.dev/living/fpga-sim/}{/living/fpga-sim/}.

\textbf{Scope of Part I.} Section I.2 places the commit record against
notation research, information theory, scaling results, LUT networks,
control certificates, and runtime enforcement. Section I.3 defines the
objects, the plants, the analytical energy model, the browser
instrument, and the boards. Section I.4 reports only quantities produced
by the software reference, by the boards, or by a cited external source.
Section I.5 answers objections that have published form. Section I.6
states what would falsify the claims and what was not measured.

\subsection{I.1. Introduction}\label{i.1.-introduction}

A language model can rank actions. Ranking is not permission. The
failure mode that matters for a plant, a shell, a payment, or a
destructive laboratory step is a commit: the action runs. Soft scores
remain useful upstream. They are the wrong type at the irreversible
branch, because a probability does not name which predicate failed, and
it does not force the actuator command to zero.

Lee (2022) states a prior claim cleanly. Notation is not decoration. A
notation that makes an operation executable changes which errors can be
written down and which can be hidden. Iverson (1980) made the same point
for array languages: notation is a tool of thought. Engelbart (1962)
treated language, artifacts, methods, and training as one system. Kay
(1972) and Victor (2013; 2012 talk) argued that representations should
keep consequences near the idea being edited. Matuschak and Nielsen
asked why tools for thought stall as products. That literature is about
cognition and media. It does not, by itself, specify a record that a
machine must obey before it moves.

The industrial center of recent artificial intelligence is scale. Kaplan
et al.~(2020) and Hoffmann et al.~(2022) report loss scaling with
parameters, data, and compute. Hooker (2021) argues that ideas spread
when they fit available hardware and kernels. Those results describe how
the capability of a proposer grows with the resources spent on it. They
do not define a commit predicate.

This paper's contribution is a software reference for that predicate,
plus a report of what the reference actually does on toy plants. The
contribution is compositional. Boolean LUT shields, energy certificates,
ternary table lookup, and tool transports each exist in prior work. The
reference wires them as one auditable decision: proposal, then
certificate, then typed refuse or commit. The claim is the existence and
behavior of this reference, not priority.

Three measurement statements constrain every later number. First, joules
computed by the OpCounter are analytical estimates. Second, Verilog
results are simulation. Third, the only board joules are the KV260
board-level differences of Section I.3.7, filed as \texttt{measured\_j}
with the meter named. The board-result flag is true for those KV260 runs
alone. The Pt V2 stays unmetered. A reader should be able to tell
modeled, simulated, and board-measured apart from the method clause
attached to the number.

Figure 1 draws the commit gate.

\includegraphics[width=0.95\linewidth,height=\textheight,keepaspectratio]{fig/commit-gate.png}

\emph{Figure 1. The commit gate. Diagram of propose, certify, commit or
refuse, and receipt, as defined in Section I.3. A
\href{https://research.openie.dev/figures/gates/commit-gate.png}{PNG
copy} is provided.}

\subsection{I.2. Related work}\label{i.2.-related-work}

\subsubsection{I.2.1 Notation and tools for
thought}\label{i.2.1-notation-and-tools-for-thought}

Lee (2022) argues that inventing notations can matter more than
inventing additional automated tools, because a notation changes the
thoughts that are cheap. The interview record with The Gradient is
consistent with that essay and is not an independent empirical result.
Iverson (1980), Turing Award lecture, is the classical computer-science
statement: notation is a tool of thought, and executability is part of
the test. Engelbart (1962) locates intelligence in a human plus artifact
system, not in a bare brain. Kay (1972) specifies a personal dynamic
medium. Victor's ``Media for Thinking the Unthinkable'' and ``Inventing
on Principle'' argue that creators need representations whose
consequences are visible while the idea is still being formed. Matuschak
and Nielsen, ``How can we develop transformative tools for thought?'',
document adoption and institutional failure modes of such tools.

This paper uses that lineage for one limited inference. If a notation is
executable, some failures become ordinary data rather than unrepresented
accidents. A typed refuse reason is an instance. It does not support the
strong Sapir-Whorf claim that language determines thought. The strong
claim is rejected in Section I.5. The weak claim, which is Iverson's, is
that an executable notation changes the cost of operations and the
representability of errors.

Chollet (2019) defines intelligence as skill-acquisition efficiency and
grounds the definition in algorithmic information theory. That
definition is about learning new skills, not about commit.

\subsubsection{I.2.2 Information, description length, and physical
cost}\label{i.2.2-information-description-length-and-physical-cost}

Shannon (1948) defines entropy and channel capacity. Shannon (1959)
defines rate-distortion: how much description is required for a stated
fidelity. Kolmogorov (1965) defines description length by program size.
Solomonoff (1964) defines inductive inference by a mixture over
programs. Chaitin (1977) develops program-size complexity. Cover and
Thomas (2006) is the textbook spine. Jaynes (1957) uses maximum entropy
as a rule for distributions under constraints. Brillouin (1956) connects
information to physical negentropy in the Maxwell-demon line.

Cross-entropy training of a language model is an information-theoretic
objective on a data distribution. Sutskever's 2023 Simons Institute
talk, which is a talk and not a paper, discusses unsupervised learning
as compression. That fact does not yield a Kolmogorov certificate for a
particular string, and it does not yield a commit bit. Statistical
compression of a training distribution is not the same object as a
runtime allow set under a plant constraint.

Landauer (1961) proves a lower bound: logically irreversible erasure
dissipates at least \texttt{kT\ ln\ 2} per bit in the ideal model.
Bennett (1973) shows that logically reversible computation can avoid
that erasure cost in the limit. Horowitz (2014) reports practical CMOS
energy at picojoule scales for operations and memory movement, far above
the Landauer bound, and treats energy rather than transistor count as
the constraint. Mead (1990) is the neuromorphic argument for analog
efficiency. Sandberg (2016) warns that a brain power near 20 watt does
not bound the energy of an engineered system trained from scratch. A
cortical partitioning preprint (arXiv:2102.06273) further warns against
treating all of that biological power as comparable compute. This paper
uses Landauer and Horowitz as bounds and practice tables. It does not
claim that any model in the reference operates near \texttt{kT\ ln\ 2}.

\subsubsection{I.2.3 Scale and hardware
fit}\label{i.2.3-scale-and-hardware-fit}

Hestness et al.~(2017), Kaplan et al.~(2020), and Hoffmann et al.~(2022)
document empirical scaling of deep learning loss. Amodei and Hernandez
(2018), an OpenAI blog post, describe rapid growth in training compute
for landmark results. The page is a vendor essay. This paper does not
depend on a specific doubling time from it. Sutton (2019) argues that
general methods which use computation have beaten systems that encode
the contents of human knowledge as fixed features. The essay is not an
argument against every constraint. A commit predicate is a constraint on
action, not a hand-built chess feature. Hooker (2021) explains why
alternatives lose when the installed base of chips and libraries rewards
dense matrix multiplication.

This paper ranks by joules to correct completion and separates proposal
cost from permission. A price is not a term in either.

\subsubsection{I.2.4 LUT networks, ternary lookup, and differentiable
logic}\label{i.2.4-lut-networks-ternary-lookup-and-differentiable-logic}

Umuroglu et al.~(2020) map sparse low-bit networks to FPGA truth tables
(LogicNets). Petersen et al.~(2022) train networks of logic gates with
differentiable relaxations (DiffLogic). Bacellar et al.~(2024) train
weightless networks of LUTs (differentiable weightless neural networks).
Ma et al.~(2024) study ternary-weight language models (BitNet b1.58).
Wei et al.~(2024) implement low-bit matrix multiplication by table
lookup on CPUs (T-MAC). The TeLLMe line (arXiv:2504.16266) is a ternary
LUT matrix-multiplication design aimed at FPGA execution. Those systems
show that tables and gates are a real compute substrate. They do not, in
the papers cited, attach a port-Hamiltonian energy veto and a typed
commit envelope to a tool executor.

The reference uses DiffLogic as a thin allow table on toy features, and
TLMM as a proposal path. That is a different job from using a logic
network as the whole classifier, which is the usual DiffLogic evaluation
(vision benchmarks, large gate counts). Section I.6 records the capacity
gap.

\subsubsection{I.2.5 Certificates, shields, and runtime
monitoring}\label{i.2.5-certificates-shields-and-runtime-monitoring}

Ames et al.~(2019) survey control barrier functions as set-invariance
certificates, typically enforced by a filter or a quadratic program.
Alshiekh et al.~(2018 preprint arXiv:1708.08611) introduce shields that
replace unsafe actions with safe ones, minimally when possible. Dawson,
Gao, and Fan (2022) survey learned Lyapunov and barrier certificates and
the failure modes of treating a neural network as a certificate without
a check. Manek and Kolter (2020) learn stable dynamics with a Lyapunov
structure. Greydanus, Dzamba, and Yosinski (2019) learn Hamiltonian
neural networks. Roth et al.~(2025) study stable port-Hamiltonian neural
networks. Yu, Zikelic, and Henzinger (2024) repair neural certificates
using runtime monitors. Leung and Paré (arXiv:2512.24493) give Bayesian
safety guarantees for port-Hamiltonian systems with learned energy
functions. Their energy is the plant's physical Hamiltonian. It is not
the energy of computation. The Safe path of this reference uses energy
in that same physical sense. The joules of computation are a separate
account (Section I.3.4). Sanchez et al.~(2018) survey
runtime-verification taxonomies.

\textbf{The gate is a Simplex.} Sha's Simplex architecture runs an
unverified high-performance controller beside a verified safety
controller, and a decision module hands control to the safety controller
before the plant leaves the region it can recover from (Sha 2001). ``A
proposal may not authorize itself'' is that rule, carried from control
loops to tool calls and actuators. Runtime assurance is its name in
avionics and robotics. Three older results fix the rest of the shape.
Saltzer and Schroeder named complete mediation: every access is checked
against authority, every time (1975). A reference monitor is the
mechanism that does the checking. Gray's two-phase commit separates
prepare from commit, so no participant commits alone (1978). Necula's
proof-carrying code ships a proof that the host checks before untrusted
code runs (1997). The commit record here is a reference monitor in front
of irreversible effects, with a Simplex switch to refusal and a
certificate the host checks, as in proof-carrying code. What the
reference adds to that lineage is one conjunct and one ledger: an energy
predicate inside the certificate, and joules on the same record as the
decision.

The reference does not replace that mathematics. The Safe path is a
conservative fixed-point test of a known energy identity on a pendulum,
\texttt{V\ =\ (1/2)(g\ theta\^{}2\ +\ omega\^{}2)} with
\texttt{Vdot\ =\ omega\ *\ u}. The CBF used in the toy bake-off is an
energy-set inequality on the same \texttt{V}, not a claim to have solved
the Ames quadratic program on a manipulator. Learned Lyapunov structure
on a cart-pole is a second toy, reported in Section I.4, and is not a
region-of-attraction theorem for arbitrary plants.

\subsubsection{I.2.6 World models, typed decisions, and tool
transport}\label{i.2.6-world-models-typed-decisions-and-tool-transport}

Ha and Schmidhuber (2018) and Hafner et al.~(DreamerV3,
arXiv:2301.04104) show that learned models can propose actions from
imagined trajectories. LeCun's JEPA essays (Meta research blog) argue
for predictive world models. LeCun et al.~(2006) is the earlier
energy-based learning tutorial. Prediction error is not identically zero
in these systems. A predictor can sit upstream of a gate. It cannot
replace the gate unless its predictions are perfect and the plant model
is the plant. This paper does not compete with those systems on
prediction benchmarks.

System One models return a typed answer and a probability for each
allowed option in one forward pass, with no free-form text. TypeSafe AI
named the category with its hosted Jev model (15 September 2026). Open
models followed, among them Laya from Convai Innovations, an encoder
with decision heads. They collapse generation cost when the option set
is typed and known. Anthropic's Model Context Protocol announcement
specifies discovery and transport for tools. Transport delivers a call.
It does not evaluate \texttt{Vdot}. In the reference's MCP adapter,
irreversible tools require a certificate before the executor runs. That
is an engineering claim about this adapter, tested by the demo in
Section I.4, not a claim about every MCP server in production.

\subsubsection{I.2.7 What the composition
adds}\label{i.2.7-what-the-composition-adds}

Shields refuse unsafe actions but are not, in the cited papers, exported
as the reference's four-field commit record with an analytical joule
account. DiffLogic and LUT networks compute functions. They are not, in
those papers, a veto in front of a separate proposer and a tool
executor. Energy-aware barrier filters are continuous. The reference's
allow bit is a table. The compositional claim is the wiring: LUT allow
and energy predicate and optional barrier, recorded with a typed reason,
applied before an irreversible effect, with analytical energy kept
distinct from board power. Section I.4 states which parts of that wiring
were executed.

\subsubsection{I.2.8 Certify / commit peers (2025 to
2026)}\label{i.2.8-certify-commit-peers-2025-to-2026}

Several recent systems enforce a decision before an irreversible tool
runs.

Shi et al., \textbf{Progent}
(arXiv:\href{https://arxiv.org/abs/2504.11703}{2504.11703}), check every
tool call against symbolic least-privilege rules and use an SMT solver
so privilege can only narrow without approval (monotonic confinement).
Chen, Kang, and Li, \textbf{ShieldAgent} (ICML 2025;
\href{https://proceedings.mlr.press/v267/chen25ae.html}{proceedings.mlr.press/v267/chen25ae.html}),
enforce verifiable safety-policy circuits over agent action trajectories
before shielding. Yu, Žikelić, and Henzinger (AAAI 2025;
\href{https://doi.org/10.1609/aaai.v39i25.34840\%5D(https://doi.org/10.1609/aaai.v39i25.34840)}{{[}doi:10.1609/aaai.v39i25.34840{]}(https://doi.org/10.1609/aaai.v39i25.34840))})
repair neural certificates via runtime monitoring. A companion line
verifies ReLU control-barrier certificates online over a lookahead
region (arXiv:\href{https://arxiv.org/abs/2507.11987}{2507.11987}).

\textbf{What Part I adds.} Those peers are policy, shield, or
certificate-repair objects. This reference ships a four-field commit
envelope (\texttt{proposal} → \texttt{certificate} →
\texttt{commit\textbar{}refuse} → receipt) with a conjunctive LUT allow
bit and an energy predicate, typed refuse reasons, and analytical joules
kept distinct from board power. Transport (MCP) remains proposal. Board
joules exist only for the KV260 runs in Section I.3.7, at board level.

\subsubsection{I.2.9 Runtime
enforcement}\label{i.2.9-runtime-enforcement}

The gate is a security automaton. Schneider showed that execution
monitors enforce exactly the safety properties
\href{https://doi.org/10.1145/353323.353382}{Schneider 2000}, and
Ligatti, Bauer and Walker extended the monitor to edit automata that
suppress or insert actions
\href{https://doi.org/10.1007/s10207-004-0046-8}{Ligatti et al.~2005}.
Falcone set out when to enforce rather than verify
\href{https://doi.org/10.1007/978-3-642-16612-9_9}{Falcone 2010}, and
Bauer, Leucker and Schallhart gave runtime verdicts three values
\href{https://doi.org/10.1145/2000799.2000800}{Bauer et al.~2011}, which
map onto allow, refuse and \texttt{settle\_refuse}. What this paper adds
to that line is the energy predicate in the certificate and the receipt.
Agent guards of 2024 to 2026 have the same shape, a check between a
proposed tool call and its effect: GuardAgent
\href{https://arxiv.org/abs/2406.09187}{arXiv:2406.09187} and AgentSpec
\href{https://arxiv.org/abs/2503.18666}{arXiv:2503.18666}. PolicyGuide
argues that a check on each action cannot steer a multi-step workflow
\href{https://arxiv.org/abs/2608.19861}{arXiv:2608.19861}; that is why
Part II puts the stop on the task predicate \(C(z)\) rather than on the
single act. Multi-step irreversible work also needs compensation, which
sagas supply \href{https://doi.org/10.1145/38713.38742}{Garcia-Molina
and Salem 1987}, and ``a proposal may not authorize itself'' is the
reference-monitor principle that capability kernels such as seL4 enforce
by proof \href{https://doi.org/10.1145/1629575.1629596}{Klein et
al.~2009}. The 0-false-allow scan of Section I.4 is a test over a grid,
not a proof; Section I.6 states that limit, and it applies to every
summary of the result.

\subsubsection{I.2.10 Reference points for the board
figure}\label{i.2.10-reference-points-for-the-board-figure}

An FPGA implementation of a function is expected to cost more energy
than an ASIC; Kuon and Rose measured the dynamic power gap
\href{https://doi.org/10.1109/TCAD.2006.884574}{Kuon and Rose 2007}, and
Horowitz gives the CMOS energy per operation that frames such figures
\href{https://doi.org/10.1109/ISSCC.2014.6757323}{Horowitz 2014}. A
vendor power estimate for the same bitstream beside the SOM reading is
part of the PL-rail protocol.

\subsection{I.3. Definitions and
methods}\label{i.3.-definitions-and-methods}

\subsubsection{I.3.1 Objects}\label{i.3.1-objects}

Fix a state \texttt{x} and a candidate action \texttt{u} in a set
\texttt{U}. A proposal is a tuple

\begin{Shaded}
\begin{Highlighting}[]
\NormalTok{Proposal = (u, x, hints)}
\end{Highlighting}
\end{Shaded}

\texttt{hints} may include analytical energy quantities. A proposal has
no interpretation as permission.

A certificate is a tuple

\begin{Shaded}
\begin{Highlighting}[]
\NormalTok{Certificate = (lut\_allow, energy\_ok, cbf\_ok, reasons, metrics)}
\end{Highlighting}
\end{Shaded}

\texttt{lut\_allow} and \texttt{energy\_ok} are Booleans.
\texttt{cbf\_ok} is a Boolean when the selected compose mode requires
it, and is absent otherwise. \texttt{reasons} is a list of refuse
reasons. \texttt{metrics} includes the analytical energy account and the
board-result flag.

Compose modes in the schema README are:

\begin{Shaded}
\begin{Highlighting}[]
\NormalTok{lut\_and\_energy        : commit \textless{}=\textgreater{} lut\_allow AND energy\_ok}
\NormalTok{lut\_and\_safe          : commit \textless{}=\textgreater{} lut\_allow AND safe\_energy\_ok}
\NormalTok{lut\_and\_safe\_and\_cbf  : commit \textless{}=\textgreater{} lut\_allow AND safe\_energy\_ok AND cbf\_ok}
\end{Highlighting}
\end{Shaded}

Ablations \texttt{energy\_only} and \texttt{lut\_only} exist so a
comparison can remove one conjunct. They are not the default for
irreversible tools.

A commit decision is

\begin{Shaded}
\begin{Highlighting}[]
\NormalTok{if commit then plant\_action = u else plant\_action = 0}
\end{Highlighting}
\end{Shaded}

with the certificate attached. The executor for an irreversible tool is
called only in the first branch.

Worked record. File
\texttt{artifacts/schemas/wca.commit.v1/examples/commit\_decision\_refuse.json}
stores a refuse in which \texttt{lut\_allow} is true and
\texttt{energy\_ok} is false. The reason code is \texttt{energy\_veto}.
The fixed-point detail string is
\texttt{Vdot\_q=1461\ \textgreater{}\ eps\_q=13}. \texttt{plant\_action}
is \texttt{{[}0.0{]}}. The record sets the board-result flag to false.
One true conjunct does not commit.

\subsubsection{I.3.2 Plants and oracles}\label{i.3.2-plants-and-oracles}

The primary plant is a pendulum. Continuous reference quantities used by
the Safe check are \texttt{g\ =\ 9.81}, \texttt{epsilon\ =\ 0.05},
\texttt{V\ =\ (1/2)(g\ theta\^{}2\ +\ omega\^{}2)}, and
\texttt{Vdot\ =\ omega\ *\ u}. \texttt{energy\_ok} in the continuous
oracle means \texttt{Vdot\ \textless{}=\ epsilon}.

The legacy discrete check is a Q8.8 fixed-point approximation,
bit-matched to the Verilog module \texttt{wca\_ph\_energy.v} in
simulation. The Safe check (\texttt{SafeFixedPointPHCertificate}) uses
Q16.16 and an exact integer residual. The reported predicate is

\begin{Shaded}
\begin{Highlighting}[]
\NormalTok{energy\_ok \textless{}=\textgreater{} 4*omega\_q*u\_q + 2*(abs(omega\_q)+abs(u\_q)) + 1}
\NormalTok{              \textless{}= floor(4*epsilon*S\^{}2)}
\end{Highlighting}
\end{Shaded}

with the rounding argument given in \texttt{artifacts/RESULTS.md}: under
round-to-nearest with error at most half a least significant bit,
\texttt{Vdot} is bounded by
\texttt{(prod\ +\ (1/2)\ abs\ +\ 1/4)\ /\ S\^{}2}. The design goal is
one-sided. If the Safe predicate allows, the continuous oracle allows.
The converse is not required. Extra refuses are a measured cost, not a
contradiction of the safety direction.

A second plant, cart-pole, is used only in the learned-Lyapunov
simulation. State is \texttt{{[}x,\ xdot,\ theta,\ thetadot{]}}, with
\texttt{g\ =\ 9.81}, cart mass 1.0, pole mass 0.1, length 0.5.
\texttt{V} is a quadratic form \texttt{xi\^{}T\ P\ xi\ /\ 2} with
\texttt{P} positive definite by a Cholesky construction plus a multiple
of the identity. The fit is stochastic gradient on a hinge of
\texttt{Vdot\ +\ alpha\ V} under a linear policy, seed 42, 400
iterations, batch 32, \texttt{alpha\ =\ 0.15}, as recorded in
\texttt{RESULTS.md}. The continuous gate is
\texttt{Vdot\ \textless{}=\ 0.05}. The Safe discrete rule adds a
Lipschitz margin times half a least significant bit. This is a
simulation of one fitted quadratic on one box. It is not a proof for
other plants.

\subsubsection{I.3.3 LUT allow and proposal
path}\label{i.3.3-lut-allow-and-proposal-path}

The allow LUT in the reported episodes is a DiffLogic export with mask
\texttt{21887} on four features: absolute theta, absolute omega,
absolute \texttt{u}, and an energy bit, unless a section says otherwise.
TLMM is a grouped ternary table lookup used as a proposer.
Iso-correctness means agreement with a naive ternary matrix product
within absolute tolerance 1e-12, which \texttt{RESULTS.md} records as
measured in software. Scale behavior is whatever \texttt{wca-tlmm-scale}
writes under a stated analytical forward budget. The default budget
discussed in results is 2e-10 joule per forward, analytical, not
electrical.

\subsubsection{I.3.4 Analytical energy
model}\label{i.3.4-analytical-energy-model}

\texttt{RESULTS.md} separates two facts. Operation counts are exact in
the OpCounter. Joules are the counts times constants:

{\def\LTcaptype{none} % do not increment counter
\begin{longtable}[]{@{}lll@{}}
\toprule\noalign{}
Constant & Value (joule) & Role \\
\midrule\noalign{}
\endhead
\bottomrule\noalign{}
\endlastfoot
\texttt{E\_LUT\_READ} & 1e-12 & LUT read \\
\texttt{E\_BRAM\_WORD} & 1e-11 & BRAM word \\
\texttt{E\_BITMASK\_OP} & 5e-14 & bitmask operation \\
\texttt{E\_ADD} & 5e-13 & add \\
\texttt{E\_PACK} & 2e-13 & pack \\
\texttt{E\_COMMIT\_BASE} & 2e-11 & commit overhead \\
\texttt{E\_REFUSE\_OVERHEAD} & 5e-12 & refuse overhead \\
\end{longtable}
}

\texttt{J\ =\ sum\ count\_op\ *\ E\_op}. These constants are engineering
stand-ins, not measurements. Structure counters such as shared BRAM hits
are reported and are not charged unless they change a counted read.
Changing a constant rescales every analytical joule in lockstep. That is
why the numbers cannot be compared to a vendor's tokens per joule, to
MLPerf Tiny energy, or to a rail measurement. They can be compared
across modes inside the same model.

Utility in the internal tables is commit count. Commit count is not a
task reward. A policy that commits more can score higher utility per
joule while violating the energy oracle. Section I.4 therefore does not
treat utility per joule as a safety metric.

\subsubsection{I.3.5 Tool adapter}\label{i.3.5-tool-adapter}

\texttt{wca-mcp-gate} classifies tools. Read-only tools may execute
without a commit record. Unknown tools are treated as irreversible.
Irreversible tools require a certificate under the configured compose
mode, default \texttt{lut\_and\_safe\_and\_cbf} in
\texttt{artifacts/mcp\_gate\_demo.json}. On refuse, \texttt{executed} is
false and \texttt{executor\_calls} is 0. The demo executor records calls
and does not perform host input or output. The method string in the demo
file is \texttt{rust\_mcp\_commit\_gate\_v1}.

\subsubsection{I.3.6 Simulation of
Verilog}\label{i.3.6-simulation-of-verilog}

The repository emits Verilog and memory images for the allow LUT and the
commit gate (\texttt{artifacts/fpga/}). \texttt{wca-rtl-verify} and
Icarus Verilog check bit-exact behavior of those models against the
software oracle. Icarus executes a simulation. It does not place, route,
or measure a chip. Vivado reports enter only for the KV260 runs of
Section I.3.7, as utilization and timing closure. Vivado power estimates
are not results.

\subsubsection{I.3.7 Browser instrument and
boards}\label{i.3.7-browser-instrument-and-boards}

The completeness path for replication has three stages. Stages A and B
exist. Stage C has two boards: one for decision agreement and one for
joules.

Stage A, present. Rust binaries in \texttt{wca-commit} run the episode,
the Safe check, the tool demo, and the analytical joule account. Icarus
simulates emitted Verilog on a workstation.

Stage B, shipped. The pure decision core (proposal in, certificate out,
plant step, analytical joule update) compiles to \texttt{wasm32}
(\texttt{crates/wca-fpga-sim}). A static page at
https://research.openie.dev/living/fpga-sim/ loads the module. WebGPU
holds the LUT words and a trace buffer of
\texttt{(step,\ lut\_allow,\ energy\_ok,\ cbf\_ok,\ decision)} for
display. WebGPU timestamps are not joules. On seed-1 (16 steps, theta
1.5, omega 3.0) the WASM episode reports committed 1, refused 15,
analytical joule 8.6795e-10, matching Stage A. The WASM module is an
emulation of the same functions the Rust crate runs. It is not a
cycle-accurate model of an Artix-7, not a switching-activity power
model, and not a substitute for Icarus on the Verilog. Stage B does not
claim an Alchitry Pt V2 DUT.

Stage C, board programmed, energy unmetered. An Alchitry Pt V2 was
detected and loaded (SRAM) with a closed-loop WCA commit-gate bitstream:
on-fabric seed-1 pendulum plant dynamics, TLMM proposal, allow LUT, Safe
Q16.16 port-Hamiltonian residual, and Ames-style energy-set control
barrier function (CBF) residual. Fabric commit is
\texttt{lut\_allow\ ∧\ energy\_ok\ ∧\ cbf\_ok}. UART at 115200 8N1
agreed with the software composite golden on allow/refuse/commit (1
commit, 15 refuses, sole commit at step 6). On that seed-1 closed-loop
episode the CBF conjunct did not change the allow/refuse trace versus
LUT and Safe PH alone (the sole PH commit has u=0, so the CBF Lie check
holds). \textbf{Closed-loop} here means plant state evolves from gate
commits and refuses on the fabric; it is not a streamed stimulus ROM of
plant/proposal state. The Pt V2 has \textbf{no onboard joule meter}. No
USB inline meter or shunt was attached. Vivado / post-PAR power
estimates are not DUT readings. SparkFun's product page names FPGA
XC7A100T-2FGG84I and lists 101,440 logic cells, 240 DSP48E1 slices,
4,860 Kb of block RAM, and 256 MB of DDR3L
(https://www.sparkfun.com/alchitry-pt-v2.html); those remain vendor
specifications. Builder metering methods:
\href{https://research.openie.dev/living/fpga-sim/alchitry/\#energy}{/living/fpga-sim/alchitry/\#energy}.

Stage C, board metered. An AMD Kria KV260 carries an onboard INA260
power monitor (\texttt{ina260\_u14}). Three rungs ran on 30 September
2026, each as interleaved A/B/A passes with N = 4, settle 15 s and
sampling 25 s at about 4 Hz. Rung 1 compares a free-running seed fabric
(439 LUTs) to the board's baseline bitstream. It gives watts only, with
no operation count, so it gives no joules per operation. Rung 2 halts
and runs one seed bitstream (975 LUTs) at about 100 MHz under an
on-device operation counter. The difference is +0.0098 W (standard error
0.0031 W), which is 98.0 pJ per seed cycle (standard error 30.8 pJ).
Rung 3 halts and runs the gate bitstream: eight lanes, each the
conjunction of the allow LUT, the Safe Q16.16 energy residual and the
barrier residual, in 6,370 LUTs and 128 DSP slices, timing closed at 25
MHz. The difference is +0.2857 W (standard error 0.0032 W). At 2.4999 ×
10⁷ gate cycles per second that is 11,428 pJ per gate cycle (standard
error 129 pJ), about 1.43 nJ per lane decision. The A2 minus A1 drift
was +0.0047 W (standard error 0.0040 W). The bitstream hashes and logs
are in the \texttt{wca-kv260-meter} artifact of the reference.

What the meter can and cannot see. The INA260 specifies a maximum offset
of 5 mA and a maximum system gain error of 0.15\%, with a current step
of 1.25 mA (TI INA260 datasheet). The offset is common to A and B on the
same board and rail, so the halt-versus-run difference cancels it. Gain
error scales the difference by at most 0.15\%. The meter senses board
power, so the difference includes clock tree and routing switched by the
enable, not the gate logic alone. Rung 2's difference is about three
standard errors from zero. Rung 3's is about ninety. The next protocol
is fixed: read the core rail (VCCINT) through the module's regulator
telemetry or a source-measure unit; at least 30 randomized interleaved
A/B pairs; at least 10⁹ gate evaluations per block; die temperature
logged with an idle baseline before each block; gross and subtracted
values both reported; the Vivado power estimate reported beside the
reading with the discrepancy. Rung 3 is not a closed-loop plant joule
cost, and the analytical constants of Section I.3.4 are not calibrated
by it.

The scientific reason for stage B is replication and inspection, not a
new physical claim. A reader with a browser steps the same seed-1
episode and reads the same allow bit and reason codes that the crate
prints. Agreement with stage A is the acceptance test for stage B, and
it passes on seed 1. Agreement of decisions is the acceptance test for
the Pt V2, and it passes on seed 1. A meter reading on a stated workload
is the acceptance test for the KV260, and rungs 2 and 3 meet it at board
level.

\subsubsection{I.3.8 What is not a method of this
paper}\label{i.3.8-what-is-not-a-method-of-this-paper}

No human-subject study. No claim about linguistic relativity beyond the
rejection in Section I.5. No training-compute extrapolation. No
identification of economic satiation. No optimization against Groq,
Etched, or Taalas throughput. No world-model benchmark.

\subsection{I.4. Results}\label{i.4.-results}

All quantities in this section are simulated in the software reference
unless a sentence cites an external paper. Analytical joules use Section
I.3.4. Only NI-9 is board-measured.

\subsubsection{I.4.1 Claim NI-1. The commit record exists as a schema
and a Rust
type}\label{i.4.1-claim-ni-1.-the-commit-record-exists-as-a-schema-and-a-rust-type}

Schemas \texttt{Proposal}, \texttt{Certificate}, \texttt{RefuseReason},
and \texttt{CommitDecision} are in
\texttt{artifacts/schemas/wca.commit.v1/}. Rust types are in
\texttt{crates/wca-commit/src/schema.rs}. Round-trip tests are
\texttt{cargo\ test\ -p\ wca-commit}. This is an existence result about
the reference, version \texttt{wca.commit.v1}. It is not a uniqueness
theorem.

\subsubsection{I.4.2 Claim NI-2. Seed-1 pendulum
episode}\label{i.4.2-claim-ni-2.-seed-1-pendulum-episode}

Command, from \texttt{RESULTS.md}:

\begin{Shaded}
\begin{Highlighting}[]
\NormalTok{cargo run {-}{-}release {-}{-}bin wca{-}soc{-}loop {-}{-} \textbackslash{}}
\NormalTok{  {-}{-}steps 16 {-}{-}init{-}theta 1.5 {-}{-}init{-}omega 3.0 {-}{-}seed 1}
\end{Highlighting}
\end{Shaded}

Reported outcome: committed = 1, refused = 15, analytical episode energy
= 8.6795e-10 joule. The same committed and refused counts and the same
analytical energy are reported for \texttt{-\/-safe-ph} on that seed.
\texttt{RESULTS.md} states that the single commit is step 6. Toolchain
note in that file: Rust 1.98. Re-running the binary is the reproduction.
The energy is not a measurement of a chip.

\subsubsection{I.4.3 Claim NI-3. False allows of the Safe energy
check}\label{i.4.3-claim-ni-3.-false-allows-of-the-safe-energy-check}

\texttt{wca-ph-tighten} compares discrete decisions to the continuous
oracle \texttt{Vdot\ \textless{}=\ epsilon}. On a 21 cubed grid plus
200,000 near-threshold Monte Carlo samples (209,261 decisions):

{\def\LTcaptype{none} % do not increment counter
\begin{longtable}[]{@{}lllll@{}}
\toprule\noalign{}
Path & False allow & False refuse & Both allow & Both refuse \\
\midrule\noalign{}
\endhead
\bottomrule\noalign{}
\endlastfoot
Legacy Q8.8 & 30836 & 105 & 104711 & 73609 \\
Safe v1 (Q8.8 over-approximation) & 0 & 42430 & 62386 & 104445 \\
Safe v2 (Q16.16 residual) & 0 & 79 & 104737 & 104445 \\
\end{longtable}
}

\texttt{RESULTS.md} also states that a denser 41 cubed grid plus 50,000
Monte Carlo samples, checked in \texttt{cargo\ test}, asserts false
allow equal to 0. Safe v2 cuts false refuses from 42,430 to 79 relative
to Safe v1 while keeping false allows at 0. Legacy Q8.8 is not safe on
this definition: it records 30,836 false allows. These counts are
properties of the numerical test, not of a physical pendulum.

\subsubsection{I.4.4 Claim NI-4. Toy barrier comparison, neither
definition
dominates}\label{i.4.4-claim-ni-4.-toy-barrier-comparison-neither-definition-dominates}

\texttt{wca-cbf-bakeoff} runs seeds 0 through 4, 40 steps, time step
0.05, action scale 2.5, shared demo TLMM and DiffLogic mask
\texttt{21887}. The energy-set barrier uses \texttt{h\ =\ 18\ -\ V} and
allows when \texttt{h\ \textgreater{}=\ 0} and
\texttt{omega\ *\ u\ \textless{}=\ kappa\ *\ h} with
\texttt{kappa\ =\ 0.05}. The discrete shield allows a step that stays
inside absolute theta at most 2 and absolute omega at most 4.
Closed-loop means differ because trajectories diverge. The file reports
per-episode means for commits and sums for violation counts. Selected
rows:

{\def\LTcaptype{none} % do not increment counter
\begin{longtable}[]{@{}
  >{\raggedright\arraybackslash}p{(\linewidth - 12\tabcolsep) * \real{0.0952}}
  >{\raggedright\arraybackslash}p{(\linewidth - 12\tabcolsep) * \real{0.0952}}
  >{\raggedright\arraybackslash}p{(\linewidth - 12\tabcolsep) * \real{0.0952}}
  >{\raggedright\arraybackslash}p{(\linewidth - 12\tabcolsep) * \real{0.1667}}
  >{\raggedright\arraybackslash}p{(\linewidth - 12\tabcolsep) * \real{0.2262}}
  >{\raggedright\arraybackslash}p{(\linewidth - 12\tabcolsep) * \real{0.1071}}
  >{\raggedright\arraybackslash}p{(\linewidth - 12\tabcolsep) * \real{0.2143}}@{}}
\toprule\noalign{}
\begin{minipage}[b]{\linewidth}\raggedright
Filter
\end{minipage} & \begin{minipage}[b]{\linewidth}\raggedright
Commit
\end{minipage} & \begin{minipage}[b]{\linewidth}\raggedright
Refuse
\end{minipage} & \begin{minipage}[b]{\linewidth}\raggedright
Analytical J
\end{minipage} & \begin{minipage}[b]{\linewidth}\raggedright
False allow vs PH
\end{minipage} & \begin{minipage}[b]{\linewidth}\raggedright
\texttt{h\ \textless{}\ 0}
\end{minipage} & \begin{minipage}[b]{\linewidth}\raggedright
\texttt{Vdot\ \textgreater{}\ epsilon}
\end{minipage} \\
\midrule\noalign{}
\endhead
\bottomrule\noalign{}
\endlastfoot
\texttt{lut\_and\_safe\_ph} & 1 & 38 & 2.153e-09 & 0 & 4 & 0 \\
\texttt{cbf\_energy\_set} & 1 & 38 & 2.056e-09 & 3 & 0 & 3 \\
\texttt{shield\_discrete} & 18 & 21 & 2.306e-09 & 90 & 38 & 90 \\
\texttt{lut\_only} & 40 & 0 & 2.544e-09 & 194 & 220 & 194 \\
\end{longtable}
}

On this plant and this oracle, \texttt{lut\_and\_safe\_ph} has zero
\texttt{Vdot} violations and four \texttt{h\ \textless{}\ 0} events. The
energy-set barrier has zero \texttt{h\ \textless{}\ 0} events and three
\texttt{Vdot} violations. The shield commits more and violates the
energy oracle often. \texttt{lut\_only} matches \texttt{always\_allow}
on violations in the reported table. Offline agreement on a shared
always-allow stream of 200 steps gives 4 disagreements between
\texttt{lut\_and\_safe\_ph} and \texttt{cbf\_energy\_set}, and 70
against the shield.

Interpretation inside the test: passivity and set invariance are
different predicates. The LUT-and-energy gate is not equivalent to an
Ames filter, and it does not win every column. A full quadratic-program
baseline on a richer plant is not in this table. That comparison remains
open. The result that is available is the disagreement, which is
evidence against collapsing the two certificates into one word.

\subsubsection{I.4.5 Claim NI-5. Cart-pole Safe Lyapunov
simulation}\label{i.4.5-claim-ni-5.-cart-pole-safe-lyapunov-simulation}

On the audit described in Section I.3.2 (grid n = 7 plus 20,000 boundary
samples inside the Lipschitz box, 26,487 decisions), Safe learned
Lyapunov versus the continuous \texttt{Vdot} oracle records false allow
0, false refuse 3, both allow 15,925, both refuse 10,559. Closed-loop
means over five seeds: \texttt{lut\_and\_safe\_lyap} commits 11, refuses
28, analytical J 2.719e-09, and records zero false allows against that
oracle. The level-set barrier on the same \texttt{V} commits 14 and
records 15 energy-oracle violations. The discrete shield commits 29 and
records 95. Same pattern as the pendulum: the Safe conjunct tracks its
own oracle; a different certificate tracks a different one;
permissiveness without the oracle is not safety.

This is one additional simulated plant. It does not establish the commit
gate on industrial dynamics.

\subsubsection{I.4.6 Claim NI-6. MCP
demo}\label{i.4.6-claim-ni-6.-mcp-demo}

\texttt{artifacts/mcp\_gate\_demo.json} records four cases. The
board-result flag in the file is false.

{\def\LTcaptype{none} % do not increment counter
\begin{longtable}[]{@{}lllll@{}}
\toprule\noalign{}
Case & Class & Executed & Executor calls & Commit \\
\midrule\noalign{}
\endhead
\bottomrule\noalign{}
\endlastfoot
\texttt{read\_only\_bypass} & read-only & true & (no commit record) &
absent \\
\texttt{irreversible\_refuse} & irreversible & false & 0 & false \\
\texttt{irreversible\_allow} & irreversible & true & 1 & true \\
\texttt{shell\_exec\_gated} & irreversible & true & (gated) & true \\
\end{longtable}
}

The refuse case lists \texttt{energy\_veto} (\texttt{vdot} 6.0 against
\texttt{eps} 0.05) and \texttt{cbf\_veto}. Plant action on that case is
\texttt{{[}0.0{]}}. The allow case uses compose
\texttt{lut\_and\_safe\_and\_cbf} and plant action \texttt{{[}-0.02{]}}.
The note in the file states that refuse never executes. This is a demo
of the adapter, not a field study of computer-use agents.

\subsubsection{I.4.7 Claim NI-7. TLMM under an analytical
budget}\label{i.4.7-claim-ni-7.-tlmm-under-an-analytical-budget}

\texttt{RESULTS.md} reports a demo 2 by 4, group 2 configuration near
4.24e-11 analytical joule, under the forward budget. The largest
configurations that remain under the 2e-10 forward budget stop at 64
parameters. Widths 8 by 16 and above exceed that budget even after the
optimizer. Iso-correctness against naive ternary arithmetic is a
software comparison at tolerance 1e-12. These statements do not say that
the design matches TeLLMe throughput or any board energy.

\subsubsection{I.4.8 Claim NI-8. DiffLogic allow versus an energy
teacher, toy
only}\label{i.4.8-claim-ni-8.-difflogic-allow-versus-an-energy-teacher-toy-only}

On the hard multi-seed gate benchmark in \texttt{RESULTS.md}, after
threshold calibration and denser labels, DiffLogic matches the energy
teacher on refuse rate 50.0 plus or minus 15.9 percent, precision at
threshold 0.890 plus or minus 0.067, and recall 0.802 plus or minus
0.076. A bitmask reference refuses more (71.0 plus or minus 14.9
percent) and recalls less (0.454 plus or minus 0.266). Before
calibration, DiffLogic refuse rate was 57.5 plus or minus 14.1 percent
with recall 0.724 plus or minus 0.163. The match is on this toy labeling
setup. It is not an ImageNet or control-suite result. Internal utility
per joule rises for the calibrated gate (3.570e9 versus 3.212e9 before)
under the analytical model. As Section I.3.4 states, that ratio uses
commit count as utility.

\subsubsection{I.4.9 Claim NI-9. Board-level joules of the gate
fabric}\label{i.4.9-claim-ni-9.-board-level-joules-of-the-gate-fabric}

{\def\LTcaptype{none} % do not increment counter
\begin{longtable}[]{@{}
  >{\raggedright\arraybackslash}p{(\linewidth - 12\tabcolsep) * \real{0.1429}}
  >{\raggedright\arraybackslash}p{(\linewidth - 12\tabcolsep) * \real{0.1429}}
  >{\raggedright\arraybackslash}p{(\linewidth - 12\tabcolsep) * \real{0.1429}}
  >{\raggedright\arraybackslash}p{(\linewidth - 12\tabcolsep) * \real{0.1429}}
  >{\raggedright\arraybackslash}p{(\linewidth - 12\tabcolsep) * \real{0.1429}}
  >{\raggedright\arraybackslash}p{(\linewidth - 12\tabcolsep) * \real{0.1429}}
  >{\raggedright\arraybackslash}p{(\linewidth - 12\tabcolsep) * \real{0.1429}}@{}}
\toprule\noalign{}
\begin{minipage}[b]{\linewidth}\raggedright
Rung
\end{minipage} & \begin{minipage}[b]{\linewidth}\raggedright
Fabric
\end{minipage} & \begin{minipage}[b]{\linewidth}\raggedright
Clock
\end{minipage} & \begin{minipage}[b]{\linewidth}\raggedright
Difference (W)
\end{minipage} & \begin{minipage}[b]{\linewidth}\raggedright
Operation
\end{minipage} & \begin{minipage}[b]{\linewidth}\raggedright
Joules per operation
\end{minipage} & \begin{minipage}[b]{\linewidth}\raggedright
Label
\end{minipage} \\
\midrule\noalign{}
\endhead
\bottomrule\noalign{}
\endlastfoot
1 & Seed freerun vs board baseline & n/a & −0.8212 (SE 0.0024) & none &
none & \texttt{measured\_j}, watts only \\
2 & Seed, halt vs run & about 100 MHz & +0.0098 (SE 0.0031) & one seed
cycle & 98.0 pJ (SE 30.8) & \texttt{measured\_j}, INA260, SOM rail
(\texttt{VCC\_SOM}) \\
3 & Gate, 8 lanes, halt vs run & 25 MHz & +0.2857 (SE 0.0032) & one gate
cycle, 8 decisions & 11,428 pJ (SE 129) & \texttt{measured\_j}, INA260,
SOM rail (\texttt{VCC\_SOM}) \\
\end{longtable}
}

N = 4 interleaved A/B/A passes per rung. Rung 1's negative sign means
the seed bitstream draws less than the board's baseline bitstream. It
does not mean the seed is free. Section I.3.7 gives the limits of a
board-level reading.

\subsubsection{I.4.10 Safety and
liveness}\label{i.4.10-safety-and-liveness}

Zero false allows is safety. A gate that refuses everything is also
safe, and useless. Liveness is the other axis: false refuses, and the
share of tasks the gated agent completes. NI-3 reports both on the Safe
grid. Safe v2 cuts false refuses from 42,430 to 79 at zero false allows.
The seed-1 episode commits 1 step of 16, and that count alone does not
say whether the plant reached its goal. The required report is a
frontier per plant: task completion rate against false-allow rate, one
point per compose mode of NI-4 (\texttt{lut\_only},
\texttt{shield\_discrete}, \texttt{cbf\_energy\_set},
\texttt{lut\_and\_safe\_ph}), each point with its seeds and spread. On
agent workloads the same frontier uses an attack suite for false allows
and a benign task suite for completion. A gate is judged on the
frontier, not on one axis.

\subsubsection{I.4.11 What the literature search
supports}\label{i.4.11-what-the-literature-search-supports}

A literature search in September 2026 (Sections I.2.4 to I.2.9) found no
published system that composes a learned Boolean allow LUT, a
port-Hamiltonian or Lyapunov energy predicate, and a ternary LUT
proposal as one commit boundary with bit-exact RTL simulation and an
analytical joule account. The nearest neighbors are shields,
energy-aware barrier filters, logic networks as the whole model, and
TLMM designs without an energy certificate. The search supports a gap
statement about the literature that was searched. It does not support a
priority claim over unpublished work, and it does not support an
energy-leadership claim.

\subsection{I.5. Discussion}\label{i.5.-discussion}

\subsubsection{I.5.1 Scale does not answer the commit
question}\label{i.5.1-scale-does-not-answer-the-commit-question}

Kaplan, Hoffmann, Hooker, and Sutton describe how loss and proposal
capability scale. A larger proposer can still emit an action that fails
\texttt{Vdot\ \textless{}=\ epsilon}. Putting that proposer behind the
certificate uses scale where the evidence says scale works, and uses the
predicate where scale is silent. The Bitter Lesson objects to frozen
human features that block learning. It does not entail deleting a safety
predicate. If a future result shows that an unconstrained policy meets
the same false-allow test at lower analytical cost and equal task error,
NI-3's engineering preference for the gate would be weakened on that
plant. The logical distinction between proposal and permission would
remain.

\subsubsection{I.5.2 Natural language and MCP are transport and
proposal}\label{i.5.2-natural-language-and-mcp-are-transport-and-proposal}

MCP solves discovery and call shape. System One solves some decisions
without free-form generation. Neither object is \texttt{energy\_ok}. The
demo in NI-6 shows a refuse that does not call the executor. A
production deployment could bypass the adapter. The claim is not that
the protocol makes bypass impossible. The claim is that permission has
to be a separate record if irreversible effects are to be auditable.

\subsubsection{I.5.3 Strong linguistic determinism is not the
claim}\label{i.5.3-strong-linguistic-determinism-is-not-the-claim}

The strong hypothesis that language determines thought is not a premise.
The operational statement is local. In this schema, a commit without a
passing certificate is not a successful record. An untyped confidence
string can be stored beside a success flag and still leave the reason
unstated. The enum removes that particular ambiguity. It does not
describe human cognition.

\subsubsection{I.5.4 Formalism can be wasteful; this predicate is
small}\label{i.5.4-formalism-can-be-wasteful-this-predicate-is-small}

Requiring a proof of the whole agent at every step would dominate the
budget. The reference formalizes one boundary. Read-only calls bypass.
The Safe residual is a fixed-point inequality, not a general SMT query.
NI-2's analytical episode energy is 8.6795e-10 joule under the stated
constants. It cannot be stated as a fraction of a generative model's
board energy, because that comparison was not measured. A split between
proposal joules and gate joules on a real agent workload is not in the
results. Until it is, the cost objection is answered for this toy gate
and remains open for large certificates.

\subsubsection{I.5.5 A wrong model with a green light is a real
failure}\label{i.5.5-a-wrong-model-with-a-green-light-is-a-real-failure}

If \texttt{V} is not the plant's energy, \texttt{energy\_ok} can be true
while the plant is not safe. The mitigation present in the reference is
a continuous twin on the toys, a separated reason code, and a refusal to
treat a probability as a certificate. The mitigation not present is
disturbance Input-to-State Stability, a verified region of attraction,
or a Bayesian credibility budget of the kind Leung and Paré study. NI-4
is the evidence that two respectable certificates disagree. That
disagreement is a reason to name the predicate, not a reason to trust
either one outside its test.

\subsubsection{I.5.6 Prior certificate theory is not duplicated
here}\label{i.5.6-prior-certificate-theory-is-not-duplicated-here}

Ames, Alshiekh, Dawson, and the runtime-monitoring papers already define
filters and monitors. NI-1's addition is the product record and the
composition with a LUT export and a tool executor in one reference.
Where the toy barrier wins on set invariance, the paper says so (NI-4,
NI-5). Composition is \texttt{cbf\_ok} as another conjunct when that
predicate is the one required. It is not a claim that the LUT computes
the Ames quadratic program.

\subsubsection{I.5.7 Predictors and specialized inference
chips}\label{i.5.7-predictors-and-specialized-inference-chips}

World models reduce some prediction error and leave a residual.
Specialized inference devices change the cost of proposals. Vendor
throughput figures are vendor reports until an independent meter repeats
them. The gate's job is the irreversible branch. If a device's only
interface is a token stream, the commit record still has to live
somewhere before actuators or irreversible tools run.

\subsubsection{I.5.8 Analytical joules are not
watts}\label{i.5.8-analytical-joules-are-not-watts}

The constants in Section I.3.4 are chosen engineering numbers. They make
counts comparable inside the repo. They are not calibrated to an Artix-7
rail. Reporting them as board watts would be a false measurement report.
The correct report is the one given: operation counts are computed;
analytical joules are modeled; board joules come only from the KV260
meter and are filed apart. Stage B of Section I.3.7 makes the same model
inspectable in a browser at
https://research.openie.dev/living/fpga-sim/. It does not change the
measurement class. The KV260 rungs do, for the fabric they ran.

\subsection{I.6. Limits and threats to
validity}\label{i.6.-limits-and-threats-to-validity}

Internal validity. Seed-1 is one initial condition. The Safe grid is
large but still a chosen box and a chosen epsilon. Monte Carlo samples
near the threshold stress the boundary the designers chose to stress. A
different epsilon or a different plant identity breaks the numerical
claim until re-run. The CBF in the bake-off is an energy-set inequality
implemented in the same crate, not an independently coded solver from an
Ames reference implementation. Disagreement is informative. Absolute
ranking against the literature's code is not established.

Construct validity. Commit count as utility does not measure task
success. False allow is defined against a named oracle. A system can be
oracle-safe and check-safe (zero false allows against
\texttt{Vdot\ \textless{}=\ epsilon}) and still fail set-safe
(\texttt{h\ \textgreater{}=\ 0}), which NI-4 shows. Readers who treat
``safe'' as one word will misread the table.

External validity. Both plants are toys. There is no manipulator,
quadrotor, or human-in-the-loop trial. The MCP demo does not call a
network. Verilog is simulated. The browser instrument is shipped as
Stage B simulation. The Alchitry Pt V2 agrees on decisions and is
unmetered. The KV260 reading is board level and open loop. Vendor
logic-cell counts are not confirmed here.

Measurement validity. Analytical constants can be edited to any scale.
Comparisons to Horowitz's picojoule tables, to Landauer's bound, or to
vendor tokens per joule are not valid with these constants. The KV260
rungs are a board-level meter under a stated workload. A core-rail
reading under the closed-loop episode is the missing measurement.
Post-place-and-route tool power would be a third class, still not a
meter, and it is not reported as a result.

Statistical reporting. Where \texttt{RESULTS.md} gives a mean and a
standard deviation, this paper copies them. It does not add confidence
intervals that were not computed. Multi-seed coverage is five seeds on
the bake-offs and the hard gate benchmark's reported spreads. That is
not a large-sample claim.

Not claimed. The following are not results: silicon joule leadership;
operation near the Landauer bound; brain-equivalent power; a win over
System One on latency; a win over world-model benchmarks; strong
linguistic determinism; identification of ``intelligence'' with commit
count.

Falsifiers. NI-2 is false if the stated command on the stated crate
revision does not print those counts. NI-3 is false if the Safe v2 test
reports a false allow on the stated grid. NI-6 is false if the refuse
case in the demo file executes. NI-1 is false if the schema and the
runtime diverge. A stage B emulator that disagrees with stage A on seed
1 fails its own acceptance test. A future meter that assigns board
energy to the analytical number without a calibration study does not
confirm the analytical model.

\subsection{I.7. Conclusion}\label{i.7.-conclusion}

The paper defined a commit decision as a conjunctive certificate over a
LUT allow bit and an energy predicate, with an optional barrier, and
showed a software reference that emits the decision before an
irreversible effect. On the reported pendulum tests, the Safe
fixed-point rule has zero false allows against its continuous oracle
\texttt{Vdot\ \textless{}=\ epsilon} on the stated test set, at the cost
of 79 false refuses in 209,261 decisions. On the reported comparisons,
that oracle is not the same as set invariance. Analytical episode energy
for the seed-1 run is 8.6795e-10 joule under published constants. Tool
refuse in the MCP demo does not call the executor. Verilog checks are
simulation. The browser WASM and WebGPU instrument is available at
https://research.openie.dev/living/fpga-sim/ so that the same decision
can be inspected without a board. It remains a simulation, not a DUT
result. The Alchitry Pt V2 agrees with the software on every decision of
seed 1 and stays unmetered. The KV260 meters the gate fabric at board
level: 11,428 pJ per eight-lane gate cycle at 25 MHz. The reference is a
measured software artifact and a board-metered fabric. The closed-loop
chip joule is the next measurement.

\subsection{I. References}\label{i.-references}

Simon, H. A., and Kadane, J. B. Optimal problem-solving search:
all-or-none solutions. Artificial Intelligence 6(3), 235-247, 1975.
https://doi.org/10.1016/0004-3702(75)90002-8

Alshiekh, M., Bloem, R., Ehlers, R., Konighofer, B., Niekum, S., and
Topcu, U. Safe reinforcement learning via shielding. arXiv:1708.08611.

Ames, A. D., Coogan, S., Egerstedt, M., Notomista, G., Sreenath, K., and
Tabuada, P. Control barrier functions: theory and applications. European
Control Conference, 2019. https://doi.org/10.23919/ECC.2019.8796030

Amodei, D., and Hernandez, D. AI and compute. OpenAI, 2018.
https://openai.com/index/ai-and-compute/

Anthropic. Introducing the Model Context Protocol.
https://www.anthropic.com/news/model-context-protocol

Bacellar, A., et al.~Differentiable weightless neural networks.
arXiv:2410.11112.

Bauer, A., Leucker, M., and Schallhart, C. Runtime verification for LTL
and TLTL. ACM Transactions on Software Engineering and Methodology
20(4), 14, 2011. https://doi.org/10.1145/2000799.2000800

Bennett, C. H. Logical reversibility of computation. IBM Journal of
Research and Development, 1973. https://doi.org/10.1147/rd.176.0525

Brillouin, L. Science and Information Theory. Academic Press, 1956.

Chaitin, G. J. Algorithmic information theory. IBM Journal of Research
and Development, 1977. https://doi.org/10.1147/rd.214.0350

Chen, Z., Kang, M., and Li, B. ShieldAgent: Shielding Agents via
Verifiable Safety Policy Reasoning. ICML 2025 (PMLR v267).
arXiv:2503.22738. https://proceedings.mlr.press/v267/chen25ae.html

Chollet, F. On the measure of intelligence. arXiv:1911.01547.

Convai Innovations. Laya model card.
https://huggingface.co/convaiinnovations/laya

Cover, T. M., and Thomas, J. A. Elements of Information Theory, 2nd
ed.~Wiley, 2006. https://doi.org/10.1002/047174882X

Dawson, C., Gao, S., and Fan, C. Safe control with learned certificates:
a survey of neural Lyapunov, barrier, and contraction methods.
arXiv:2202.11762.

Engelbart, D. C. Augmenting human intellect: a conceptual framework.
1962. https://www.dougengelbart.org/pubs/augment-3906-Framework.html

Falcone, Y. You should better enforce than verify. Runtime Verification
(RV 2010), Lecture Notes in Computer Science 6418, 89-105. Springer,
2010. https://doi.org/10.1007/978-3-642-16612-9\_9

Garcia-Molina, H., and Salem, K. Sagas. Proceedings of ACM SIGMOD 1987,
249-259. https://doi.org/10.1145/38713.38742

Gray, J. Notes on data base operating systems. In Operating Systems: An
Advanced Course, Lecture Notes in Computer Science 60, 393-481.
Springer, 1978. https://doi.org/10.1007/3-540-08755-9\_9

Greydanus, S., Dzamba, M., and Yosinski, J. Hamiltonian neural networks.
arXiv:1906.01563.

Ha, D., and Schmidhuber, J. World models. arXiv:1803.10122.

Hafner, D., Pasukonis, J., Ba, J., and Lillicrap, T. Mastering diverse
domains through world models. arXiv:2301.04104.

Henzinger, T. A., Kueffner, K., and Yu, E. Formal verification of neural
certificates done dynamically. arXiv:2507.11987.
https://arxiv.org/abs/2507.11987

Hestness, J., et al.~Deep learning scaling is predictable, empirically.
arXiv:1712.00409.

Hoffmann, J., et al.~Training compute-optimal large language models.
arXiv:2203.15556.

Hooker, S. The hardware lottery. arXiv:2009.06489. Communications of the
ACM, 2021. https://doi.org/10.1145/3467017

Horowitz, M. Computing's energy problem (and what we can do about it).
IEEE International Solid-State Circuits Conference, 2014.
https://doi.org/10.1109/ISSCC.2014.6757323

Iverson, K. E. Notation as a tool of thought. Communications of the ACM,
1980. https://dl.acm.org/doi/10.1145/358896.358899

Jaynes, E. T. Information theory and statistical mechanics. Physical
Review 106, 620, 1957. https://doi.org/10.1103/PhysRev.106.620

Kang, S., Yu, T., and Hwang, S. J. PolicyGuide: from guarding one action
to guiding the whole workflow for policy-compliant LLM agents.
arXiv:2608.19861. https://arxiv.org/abs/2608.19861

Kaplan, J., et al.~Scaling laws for neural language models.
arXiv:2001.08361.

Kay, A. A personal computer for children of all ages. 1972.
https://worrydream.com/refs/Kay\_1972\_-\_A\_Personal\_Computer\_for\_Children\_of\_All\_Ages.pdf

Klein, G., et al.~seL4: formal verification of an OS kernel. Proceedings
of SOSP 2009, 207-220. https://doi.org/10.1145/1629575.1629596

Kolmogorov, A. N. Three approaches to the quantitative definition of
information. 1965. English: International Journal of Computer
Mathematics, 1968. https://doi.org/10.1080/00207166808803030

Kuon, I., and Rose, J. Measuring the gap between FPGAs and ASICs. IEEE
Transactions on Computer-Aided Design 26(2), 203-215, 2007.
https://doi.org/10.1109/TCAD.2006.884574

Landauer, R. Irreversibility and heat generation in the computing
process. IBM Journal of Research and Development, 1961.
https://doi.org/10.1147/rd.53.0183

Laya. System One models and typed decision engines (comparison table,
checked 2 October 2026). https://laya-ai.com/system-one-models

LeCun, Y. Path towards autonomous machine intelligence / JEPA
discussion. Meta AI blog.
https://ai.meta.com/blog/yann-lecun-advances-in-ai-research/

LeCun, Y., Chopra, S., Hadsell, R., Ranzato, M., and Huang, F. J. A
tutorial on energy-based learning. 2006.
http://yann.lecun.com/exdb/publis/pdf/lecun-06.pdf

Lee, L. Notational intelligence. 2022.
https://thesephist.com/posts/notation/

Leung, C. H., and Paré, P. E. Bayesian safety guarantees for
port-Hamiltonian systems with learned energy functions.
arXiv:2512.24493.

Levy, W. B., and Calvert, V. G. Cerebral cortical communication
overshadows computational energy-use, but these combine to predict
synapse number. arXiv:2102.06273.

Ligatti, J., Bauer, L., and Walker, D. Edit automata: enforcement
mechanisms for run-time security policies. International Journal of
Information Security 4, 2-16, 2005.
https://doi.org/10.1007/s10207-004-0046-8

Ma, S., et al.~The era of 1-bit LLMs: all large language models are in
1.58 bits. arXiv:2402.17764.

Manek, G., and Kolter, J. Z. Learning stable deep dynamics models.
arXiv:2001.06116.

Matuschak, A., and Nielsen, M. How can we develop transformative tools
for thought? https://numinous.productions/ttft/

Mead, C. Neuromorphic electronic systems. Proceedings of the IEEE, 1990.
https://doi.org/10.1109/5.58356

Necula, G. C. Proof-carrying code. Proceedings of the 24th ACM
SIGPLAN-SIGACT Symposium on Principles of Programming Languages (POPL
'97), 106-119, 1997. https://doi.org/10.1145/263699.263712

Petersen, F., et al.~Deep differentiable logic gate networks.
arXiv:2210.08277.

Qiao, Y., Chen, Z., Zhang, Y., Wang, Y., et al.~TeLLMe: an
energy-efficient ternary LLM accelerator for prefilling and decoding on
edge FPGAs. arXiv:2504.16266.

Roth, F., et al.~Stable port-Hamiltonian neural networks.
arXiv:2502.02480.

Saltzer, J. H., and Schroeder, M. D. The protection of information in
computer systems. Proceedings of the IEEE 63(9), 1278-1308, 1975.
https://doi.org/10.1109/PROC.1975.9939

Sanchez, C., et al.~A survey of challenges for runtime verification from
advanced application domains. arXiv:1811.06740.

Sandberg, A. Energetics of the brain and AI. arXiv:1602.04019.

Schneider, F. B. Enforceable security policies. ACM Transactions on
Information and System Security 3(1), 30-50, 2000.
https://doi.org/10.1145/353323.353382

Sha, L. Using simplicity to control complexity. IEEE Software 18(4),
20-28, 2001. https://doi.org/10.1109/MS.2001.936213

Shannon, C. E. A mathematical theory of communication. Bell System
Technical Journal, 1948.
https://doi.org/10.1002/j.1538-7305.1948.tb01338.x

Shannon, C. E. Coding theorems for a discrete source with a fidelity
criterion. IRE National Convention Record, 1959.

Shi, T., He, J., Wang, Z., et al.~Progent: Securing AI Agents with
Privilege Control. arXiv:2504.11703. https://arxiv.org/abs/2504.11703

Solomonoff, R. J. A formal theory of inductive inference, parts I and
II. Information and Control, 1964.

SparkFun. Alchitry Pt V2 product page.
https://www.sparkfun.com/alchitry-pt-v2.html

Sutskever, I. An observation on generalization. Simons Institute talk,
2023. https://www.youtube.com/watch?v=AKMuA\_TVz3A

Sutton, R. The bitter lesson. 13 March 2019.
http://www.incompleteideas.net/IncIdeas/BitterLesson.html

Texas Instruments. INA260 precision digital current and power monitor,
data sheet. https://www.ti.com/lit/ds/symlink/ina260.pdf

TypeSafe AI. Introducing System One Models and Jev. 15 September 2026.
https://typesafe.ai/blog/introducing-system-one-models-and-jev

Umuroglu, Y., et al.~LogicNets. arXiv:2004.03021.

Victor, B. Inventing on principle. 2012.
https://www.youtube.com/watch?v=PUv66718DII

Victor, B. Media for thinking the unthinkable.
https://worrydream.com/MediaForThinkingTheUnthinkable/

Wang, H., Poskitt, C. M., and Sun, J. AgentSpec: customizable runtime
enforcement for safe and reliable LLM agents. arXiv:2503.18666.
https://arxiv.org/abs/2503.18666

Wei, J., et al.~T-MAC. arXiv:2407.00088.

Xiang, Z., Zheng, L., Li, Y., Hong, J., et al.~GuardAgent: safeguard LLM
agents by a guard agent via knowledge-enabled reasoning.
arXiv:2406.09187. https://arxiv.org/abs/2406.09187

Yu, E., Žikelić, Đ., and Henzinger, T. A. Neural control and certificate
repair via runtime monitoring. AAAI 2025.
https://doi.org/10.1609/aaai.v39i25.34840 (arXiv:2412.12996).

Local measurement record: \texttt{artifacts/RESULTS.md} in the
\texttt{wca-lut-edge} software reference. KV260 meter record:
\texttt{artifacts/stage-c/wca-kv260-meter/} (protocol, claim status,
A/B/A logs and summaries). Schema examples:
\texttt{artifacts/schemas/wca.commit.v1/examples/}. Demo:
\texttt{artifacts/mcp\_gate\_demo.json}.

\subsection{I. Appendix A.
Reproducibility}\label{i.-appendix-a.-reproducibility}

Workstation path used to generate the cited logs: the
\texttt{wca-lut-edge} tree, Rust 1.98 as recorded in
\texttt{RESULTS.md}. Day-to-day commands:

\begin{Shaded}
\begin{Highlighting}[]
\NormalTok{cargo test {-}p wca{-}commit}
\NormalTok{cargo run {-}{-}release {-}{-}bin wca{-}soc{-}loop {-}{-} \textbackslash{}}
\NormalTok{  {-}{-}steps 16 {-}{-}init{-}theta 1.5 {-}{-}init{-}omega 3.0 {-}{-}seed 1}
\NormalTok{cargo run {-}{-}release {-}{-}bin wca{-}soc{-}loop {-}{-} \textbackslash{}}
\NormalTok{  {-}{-}steps 16 {-}{-}init{-}theta 1.5 {-}{-}init{-}omega 3.0 {-}{-}seed 1 {-}{-}safe{-}ph}
\NormalTok{cargo run {-}{-}release {-}{-}bin wca{-}ph{-}tighten {-}{-} {-}{-}grid 21 {-}{-}out artifacts/ph\_tighten\_report.json}
\NormalTok{cargo run {-}{-}release {-}{-}bin wca{-}cbf{-}bakeoff {-}{-} {-}{-}out artifacts/cbf\_bakeoff\_report.json}
\NormalTok{cargo run {-}{-}release {-}{-}bin wca{-}lyap{-}plant {-}{-} {-}{-}out artifacts/lyap\_plant\_report.json}
\NormalTok{cargo run {-}{-}release {-}{-}bin wca{-}mcp{-}gate {-}{-} {-}{-}demo}
\NormalTok{cargo run {-}{-}release {-}{-}bin wca{-}tlmm{-}scale {-}{-} {-}{-}out artifacts/tlmm\_scale\_report.json}
\end{Highlighting}
\end{Shaded}

Icarus Verilog was version 12.0 in the environment note in
\texttt{RESULTS.md}. RTL checks are simulation. Do not read analytical
joules as board power.

The board-result flag is true for the KV260 Stage C rungs 1 to 3 only
(Section I.3.7), at board level. Every other result in this paper is
simulation or analytical.

\subsection{I. Appendix B. Browser FPGA emulator:
architecture}\label{i.-appendix-b.-browser-fpga-emulator-architecture}

This appendix describes the architecture of the Stage B emulator, which
runs at https://research.openie.dev/living/fpga-sim/. Seed-1 agreement
numbers below are from the WASM run against the Stage A golden.

Modules.

\begin{enumerate}
\def\labelenumi{\arabic{enumi}.}
\tightlist
\item
  \texttt{decision\_core}. Pure functions already represented by the
  Rust certificate, plant step, and joule accumulator. No file system
  and no threads required. Target \texttt{wasm32-unknown-unknown}.
\item
  \texttt{lut\_image}. The same memory image the Verilog path loads
  (\texttt{artifacts/fpga/wca\_allow\_lut.mem} and successors). The
  browser copy must hash-match the file used for the Icarus check.
\item
  \texttt{trace}. A fixed schema: step index, state, \texttt{u},
  \texttt{lut\_allow}, \texttt{energy\_ok}, \texttt{cbf\_ok}, reason
  codes, analytical joule increment, plant action. Identical field names
  to \texttt{wca.commit.v1} so a dump can be diffed against the crate's
  JSON.
\item
  \texttt{webgpu\_view}. Device, queue, and a storage buffer for the LUT
  and the trace. A shader draws allow and refuse. A second dispatch may
  map the Safe residual across a grid. Buffers are for inspection. They
  are not a power model.
\item
  \texttt{workstation\_oracle}. Icarus remains the RTL oracle. The
  browser does not replace it. Disagreement between WASM and the Rust
  binary fails the build. Disagreement between WASM and Icarus on the
  shared LUT image fails the build. Neither comparison mentions watts.
\end{enumerate}

Acceptance tests. The seed-1 test passes (Section I.3.7). The others are
the specified checks for the module.

\begin{itemize}
\tightlist
\item
  Seed-1, 16 steps, init theta 1.5, omega 3.0: committed 1, refused 15,
  analytical joule 8.6795e-10, matching NI-2.
\item
  Safe grid smoke: false allow 0 on a published subset of the
  209,261-decision test, matching NI-3's direction.
\item
  MCP refuse vector: executor call count 0, matching NI-6.
\item
  A unit test that WebGPU buffer readback equals the WASM trace. If
  readback is unavailable, the page must show that the GPU path did not
  confirm the trace, and the CPU WASM trace remains the result.
\end{itemize}

Physical follow-on for board \textbf{joules}. Stage C already programmed
a closed-loop commit-gate on the Alchitry Pt V2 with UART decision
agreement; the board still has no onboard joule meter. A USB inline
power meter or a shunt-plus-scope on the 5 V rail, on a stated workload,
is what would support a board-measured joule. Vivado / post-PAR
estimates are not that meter. Vendor logic-cell and memory figures stay
vendor figures. See
\href{https://research.openie.dev/living/fpga-sim/alchitry/\#energy}{/living/fpga-sim/alchitry/\#energy}.

Out of scope for the emulator: training DiffLogic in the browser and
estimating Artix-7 dynamic power from a shader.

\subsection{Part II. The cascade (Mixture of
Limits)}\label{part-ii.-the-cascade-mixture-of-limits}

\subsection{II.1. The rule}\label{ii.1.-the-rule}

A task is given by a completeness predicate \(C(z) \in \{0, 1\}\) over
its state \(z\). A \textbf{gear} is a mechanism that may close the task:
Lookup (read a stored, verified value), Formula (evaluate a closed
form), Solver (run a search, optimizer or settle with a certificate), or
Model (run a learned generator). Gear \(i\) has a pass probability
\(p_i\), the probability that it reaches \(C(z) = 1\) on the task class,
and an expected energy \(c_i\) in joules (J).

The \textbf{floor} of a task is the least expected energy at which some
policy over gears reaches \(C(z) = 1\) with failure probability at most
\(\delta\). It equals \(J^*\) of Section 0.2. The \textbf{plateau} is
the region of spend beyond the floor: once \(C(z) = 1\), no further
computation can raise \(C\), so further joules buy nothing on the stated
objective.

\textbf{Mixture of Limits} is the rule: try gears in an order that
minimizes expected energy to the first close, stop at the first gear
that closes \(C(z)\), and refuse with a receipt when the next gear would
exceed the budget or cannot change the decision. Lookup → Formula →
Solver → Model is the standard order. It is optimal only under the
condition derived in Section II.2. Every close passes through Part I's
record: propose, certify, commit or refuse, receipt.

\textbf{Where the model sits.} Model last is an order of trial, not a
verdict on what models can do. Snell et al.~(2024) showed that a
compute-optimal schedule of test-time compute beats best-of-N by more
than 4× in efficiency, and that on some problems test-time compute can
substitute for a larger model. That result concerns the schedule inside
the Model gear, which is priced in joules like any other gear. The Model
gear opens whenever cheaper gears cannot close \(C(z)\), and it opens
first when they pass too rarely to pay for their failures.

Figure 2 draws the floor and the plateau for one task.

\includegraphics[width=0.95\linewidth,height=\textheight,keepaspectratio]{fig/mol-plateau.png}

\emph{Figure 2. Mixture of Limits on one task. Diagram, not a
measurement. A
\href{https://research.openie.dev/figures/gates/mol-plateau.png}{PNG
copy} is provided.}

\textbf{Mixture of Limits is not Mixture of Experts.} Mixture-of-Experts
(MoE) routing selects parameters inside one generative model (Shazeer et
al., 2017; Fedus et al., 2022). Mixture of Limits selects whether and
where to spend across gears, and an MoE model can serve as its Model
gear.

{\def\LTcaptype{none} % do not increment counter
\begin{longtable}[]{@{}
  >{\raggedright\arraybackslash}p{(\linewidth - 4\tabcolsep) * \real{0.3333}}
  >{\raggedright\arraybackslash}p{(\linewidth - 4\tabcolsep) * \real{0.3333}}
  >{\raggedright\arraybackslash}p{(\linewidth - 4\tabcolsep) * \real{0.3333}}@{}}
\toprule\noalign{}
\begin{minipage}[b]{\linewidth}\raggedright
\end{minipage} & \begin{minipage}[b]{\linewidth}\raggedright
MoE
\end{minipage} & \begin{minipage}[b]{\linewidth}\raggedright
Mixture of Limits
\end{minipage} \\
\midrule\noalign{}
\endhead
\bottomrule\noalign{}
\endlastfoot
Where the choice is made & Inside one model & Across gears, before any
model runs \\
Object chosen & Which parameters run & Which gear closes, or refuse \\
Success & Next-token likelihood & \(C(z) = 1\), a certified commit, and
a receipt \\
Refusal & Not a native outcome & A typed outcome with a receipt \\
\end{longtable}
}

\subsection{II.2. Expected energy of a cascade and the optimal
order}\label{ii.2.-expected-energy-of-a-cascade-and-the-optimal-order}

\textbf{Two gears.} Let a light gear L with pass probability \(q_L\) and
energy \(c_L\) precede a heavy gear H that always passes, with energy
\(c_H\). Trying L first costs \(c_L + (1 - q_L)\,c_H\) in expectation.
Trying H first costs \(c_H\). So L first is cheaper if and only if

\[c_L + (1 - q_L)\,c_H < c_H \iff c_L < q_L\, c_H .\]

For example, a lookup costing \(c_L = 10^{-9}\) J that passes on
\(q_L = 0.3\) of tasks ahead of a model call costing \(c_H = 10\) J
saves \(q_L c_H - c_L \approx 3\) J per task in expectation. If \(q_L\)
falls below \(c_L / c_H = 10^{-10}\), the order reverses.

\textbf{Many gears.} For \(n\) gears tried in sequence until the first
pass, with independent outcomes, the expected energy of order \(\sigma\)
is

\[\mathbb{E}[J_\sigma] = \sum_{k=1}^{n} c_{\sigma(k)} \prod_{j<k} \bigl(1 - p_{\sigma(j)}\bigr).\]

Swapping two adjacent gears \(a\) then \(b\) changes this sum by a
factor common to both terms times
\((c_a + (1-p_a)c_b) - (c_b + (1-p_b)c_a) = c_a p_b - c_b p_a\). So
\(a\) should precede \(b\) if and only if \(p_a / c_a > p_b / c_b\), and
the optimal order sorts gears by decreasing \(p_i / c_i\) (Simon and
Kadane, 1975). Cheapest first is optimal exactly when cheaper gears also
have the higher pass probability per joule. When that fails, the cascade
starts at a heavier gear. Under a binding budget the joules of failed
light attempts stay spent, so a selector that estimates \(p_i / c_i\)
per task can close tasks that a fixed order cannot.

\textbf{Physical floor.} Any gear that erases \(n\) bits irreversibly at
temperature \(T\) dissipates at least

\[E_{\text{Landauer}}(n, T) = n\, k_B T \ln 2, \qquad k_B = 1.380649 \times 10^{-23}\ \text{J/K},\]

which is \(2.87 \times 10^{-21}\) J per bit at 300 K (Landauer, 1961;
Bennett, 1973, for the reversible limit). Measured CMOS operations cost
about 0.9 pJ for a 32-bit floating-point add and about 640 pJ for a
32-bit DRAM read at 45 nm (Horowitz, 2014), eight to eleven orders of
magnitude above that floor. The Landauer figure on a receipt is
therefore a labeled lower bound, never a measurement.

\textbf{Catalog cost model.} Before a gear runs, the software reference
prices it as

\[E_{\text{est}}(x) = \theta(D)\,\mu(S, V),\]

where \(\theta(D)\) (dimensionless) scales with the difficulty class
\(D\) of act \(x\) and \(\mu(S, V)\) (J) is a catalog energy for
primitive \(S\) on hardware view \(V\). Both are stated catalog
constants, so \(E_{\text{est}}\) is an estimate used to order gears. It
is not derived from physics, and it is never reported as a measurement.
Ranking among systems that close the task uses metered joules (Section
0.2).

\subsection{II.3. Stopping conditions}\label{ii.3.-stopping-conditions}

\textbf{Value of information.} Let \(B\) be the payoff of a decision,
\(x\) the current information state, \(y\) a further observation or
computation, \(c(y)\) its energy, and \(\lambda\) the exchange rate
between payoff and joules. The net value is

\[\Delta(x, y) = \mathbb{E}[B \mid x, y] - \mathbb{E}[B \mid x] - \lambda\, c(y)\]

(Howard, 1966; Russell and Wefald, 1991). The rule refuses escalation
when \(\Delta \le 0\). Once \(C(z) = 1\), the first two terms are equal
for any further \(y\) on the task, so \(\Delta = -\lambda c(y) < 0\) and
the stop follows.

\textbf{Grammar coverage.} Let \(G\) be the set of verified derivations
available to Lookup, Formula and Solver. A query \(q\) is covered when a
derivation \(q \vdash_G a\) exists. When \(q\) is covered, the Model
gear is not opened, because a covered gear passes with probability 1 at
lower energy than the Model gear, so by the two-gear inequality of
Section II.2 it goes first and closes the task. When \(q\) is not
covered and policy forbids the Model gear, the rule refuses.

\textbf{Settle refusal.} A Solver gear that searches for a low-energy
configuration either reaches a certified fixed point (commit) or reports
that it will not settle within its budget (refuse).

Named refusal reasons in the software reference:

{\def\LTcaptype{none} % do not increment counter
\begin{longtable}[]{@{}ll@{}}
\toprule\noalign{}
Reason & Condition \\
\midrule\noalign{}
\endhead
\bottomrule\noalign{}
\endlastfoot
\texttt{voi} & \(\Delta \le 0\) for every available next step \\
\texttt{grammar} & Not covered and the Model gear is not allowed \\
\texttt{energy} & The next gear's estimate exceeds the remaining
budget \\
\texttt{efa\_certificate} & The certificate check fails before commit \\
\texttt{settle\_refuse} & The Solver does not settle within its
budget \\
\texttt{primitive\_gap} & A required primitive is absent from the
catalog \\
\texttt{safety} & A capability check denies the action \\
\end{longtable}
}

\subsection{II.4. Provenance and
receipts}\label{ii.4.-provenance-and-receipts}

Every answer carries a replay class stating how it was produced:
\texttt{Deterministic} (Lookup, Formula, or a certified Solver),
\texttt{RetrievedCited}, \texttt{Composed}, or \texttt{ModelGenerated}.
The type system forbids strengthening a class: \texttt{ModelGenerated}
cannot be rewritten as \texttt{Deterministic}. Energy on a receipt
carries one of three labels: \textbf{Estimated} (catalog, Landauer or
instruction-fuel figure), \textbf{Unmetered} (no reading), or
\textbf{Metered} (a named meter returned a reading). The software
reference leaves \texttt{measured\_j} unset unless a meter returns. The
catalog of primitives the reference navigates, grouped into families, is
called the Periodic Stack in this paper; each cell is Present
(implemented) or Gap (absent).

\subsection{II.5. Software reference}\label{ii.5.-software-reference}

The software reference \texttt{mixture-of-limits} (Apache-2.0 or MIT)
has a harness, \texttt{cargo\ run\ -p\ mol-cli\ -\/-\ prove}, that
prints \texttt{VERIFIED} for each criterion and exits 0. The workspace
test suite (\texttt{cargo\ test\ -\/-workspace}) passes (P0), and the
harness exits 0 (P1). A criterion marked verified means that an
implementation exists and passes its stated specification. It is not a
theorem, not an independent replication and not an energy measurement.

{\def\LTcaptype{none} % do not increment counter
\begin{longtable}[]{@{}
  >{\raggedright\arraybackslash}p{(\linewidth - 10\tabcolsep) * \real{0.1667}}
  >{\raggedright\arraybackslash}p{(\linewidth - 10\tabcolsep) * \real{0.1667}}
  >{\raggedright\arraybackslash}p{(\linewidth - 10\tabcolsep) * \real{0.1667}}
  >{\raggedright\arraybackslash}p{(\linewidth - 10\tabcolsep) * \real{0.1667}}
  >{\raggedright\arraybackslash}p{(\linewidth - 10\tabcolsep) * \real{0.1667}}
  >{\raggedright\arraybackslash}p{(\linewidth - 10\tabcolsep) * \real{0.1667}}@{}}
\toprule\noalign{}
\begin{minipage}[b]{\linewidth}\raggedright
Prove \#
\end{minipage} & \begin{minipage}[b]{\linewidth}\raggedright
Criterion
\end{minipage} & \begin{minipage}[b]{\linewidth}\raggedright
Claim id
\end{minipage} & \begin{minipage}[b]{\linewidth}\raggedright
Claim
\end{minipage} & \begin{minipage}[b]{\linewidth}\raggedright
Scope note
\end{minipage} & \begin{minipage}[b]{\linewidth}\raggedright
Measurement
\end{minipage} \\
\midrule\noalign{}
\endhead
\bottomrule\noalign{}
\endlastfoot
Prove \# & Criterion (PLAN) & Paper claim id & Claim (one sentence) &
Software reference / scope note & Measurement \\
--- & --- & --- & --- & --- & --- \\
P2 & Deterministic close (double-run) &
\texttt{claim.mol.deterministic\_close} & Same asks yield the same
commit/refuse + limit id. & none & Unmetered \\
P3 & Formula/lookup close without model &
\texttt{claim.mol.formula\_lookup\_no\_model} & Landauer / unit-convert
closes COMMIT at Lookup or Formula; model never answered. & none &
Unmetered / Estimated (Landauer floor on receipt) \\
P4 & VoI refuse when \texttt{!allow\_model} &
\texttt{claim.mol.voi\_refuse} & Free-form ask → REFUSE \texttt{voi}. &
none & Unmetered \\
P5 & Settle commit + settle refuse & \texttt{claim.mol.settle\_ternary}
& Ternary settle → COMMIT; \texttt{will\ not\ settle} → REFUSE
\texttt{settle\_refuse}. & Software-ref settle stub; no VM / FPGA &
Unmetered \\
P6 & Certificate refuse on diverge &
\texttt{claim.mol.efa\_certificate\_refuse} & Formula + \texttt{diverge}
→ REFUSE \texttt{efa\_certificate}. & In-tree certificate & Unmetered \\
P7 & Capability default-deny mutate &
\texttt{claim.mol.capability\_deny\_mutate} &
\texttt{AutomateGate::default().gate(Mutate)} → Refuse. & none &
Unmetered \\
P8 & Receipt labels & \texttt{claim.mol.receipt\_honesty} &
Software-ref: \texttt{measured\_j=None}, estimate labeled
(\texttt{EstimateKind}). & Live RAPL/NVML default close OUT OF SCOPE &
Unmetered / Estimated \\
P9 & \texttt{ModelGenerated} ↛ \texttt{Deterministic} &
\texttt{claim.mol.replay\_no\_strengthen} &
\texttt{TypedAnswer::weaken\_to} returns \texttt{ReplayCoercion}. & none
& none \\
P10 & Periodic Stack live catalog navigation &
\texttt{claim.mol.stack\_live\_catalog} & Family + present + scale →
COMMIT at Lookup; 258 live Lookup/Formula/Solver/Navigate gears. & μ
calib + silicon HW OUT OF SCOPE; software reference Gap sims in proof &
Unmetered \\
P11 & \texttt{primitive\_gap} via registry probe &
\texttt{claim.mol.primitive\_gap} & Gap marker \texttt{physical\_settle}
+ absent name → REFUSE \texttt{primitive\_gap} (not string-only). & none
& Unmetered \\
P12 & μ / impedance catalog & \texttt{claim.mol.mu\_catalog} & Receipts
stamp \texttt{mu\_source=catalog}, \texttt{mu},
\texttt{landauer\_floor\_ratio}; \texttt{E≈θ·μ} (catalog estimate, not
RAPL/\texttt{measured\_j}). & μ \textbf{calib corpus} OUT OF SCOPE &
Estimated \\
P13 & Receipt transcript replay & \texttt{claim.mol.receipt\_replay} &
JSONL/in-memory replay reproduces commit/refuse + limit id; model never
answered. & none & Unmetered \\
P14 & Z2 retrieve+cite & \texttt{claim.mol.z2\_retrieve\_cite} & Seeded
factual hit → COMMIT \texttt{RetrievedCited} + citation ids; unknown →
REFUSE \texttt{claim\_unknown} (no invent). & Seed corpus in-tree only &
Unmetered \\
P15 & Z1 compose/synthesis & \texttt{claim.mol.z1\_compose} & ≥2 cited
claims → COMMIT \texttt{Composed} + \texttt{synthesis.composed\_from};
missing → REFUSE \texttt{compose\_missing}; never laundered as
Deterministic/RetrievedCited alone. & none & Unmetered \\
P16 & Agent mailbox loop & \texttt{claim.mol.agent\_mailbox} &
Goal/Message/Act → close → commit\textbar refuse + receipt in
transcript; demo hits formula/cite/compose COMMIT and voi/settle\_refuse
REFUSE; model cold; measurement flags hold. & none & Unmetered \\
P17 & Bitemporal state + memory & \texttt{claim.mol.bitemporal\_memory}
& Writes only on Mixture of Limits close COMMIT; Mutate deny / free-form
remember → REFUSE; recall → \texttt{RetrievedCited} + \texttt{memory:}
cite; unknown → \texttt{memory\_unknown}. & none & Unmetered \\
P18 & Fabric routing + software reference inventory &
\texttt{claim.mol.fabric\_soft\_ref} & Offline inventory valid
(\texttt{source=software\_ref}, Cpu always); Lookup/Formula→Cpu;
Settle→ThermoSettle\textbar Cpu; Model→Gpu* or REFUSE
\texttt{fabric\_unavailable}; detect ≠ measured joules. & Live wgpu
presence optional / not required for prove & Unmetered \\
P19 & Desktop shell headless & \texttt{claim.mol.desktop\_headless} &
\texttt{ShellSession} ask/close → receipt view
(zone/fabric/limit/estimated\_j/\texttt{measured\_j=None}); joule
ledger; VoI refuse; GUI not required. & egui GUI optional & Unmetered /
Estimated \\
P20 & OS meter labels & \texttt{claim.mol.os\_meter\_honesty} & Feature
off / VM / no permission → \texttt{measured\_j=None}, empty components,
\texttt{measure\_source=unavailable}; fixtures stamp labeled sources;
default close stays unmetered. & Live probe optional
(\texttt{energy-meter}) & Unmetered (default) / Metered only when probe
returns \\
P21 & WASM capsule certify & \texttt{claim.mol.wasm\_capsule} &
\texttt{add.wasm} under stub runtime; fuel → estimated\_j only; FS/net
grant / fuel exceed / host-share → REFUSE; close stamps capsule. &
Optional \texttt{wasmtime} feature & Unmetered / Estimated (fuel) \\
P22 & Agent Lane session path & \texttt{claim.mol.agent\_lane} & Lane
partitions; \texttt{host\_invoke} without provenance refuses; keyword
confirm insufficient; need \texttt{GrantReceipt}; allow stamps
\texttt{AgentLaneReceipt}. & none & Unmetered \\
P23 & Multi-fabric compute receipt &
\texttt{claim.mol.fabric\_compute\_receipt} & Software reference
inventory; \texttt{ComputeFabric} / \texttt{ScheduleDecision} /
\texttt{ComputeStepReceipt} with \texttt{fabric\_id}; cascade stamps
\texttt{fabric:route}; CPU Formula COMMIT; Model residual may
\texttt{fabric\_unavailable}; \texttt{measured\_j=None} on software
reference path. & none & Unmetered \\
P23b & Live fabric soft & \texttt{claim.mol.fabric\_live\_soft} &
Feature off stays software\_ref; mock/live Metal can stamp
\texttt{gpu\_metal}; Formula stays Cpu; missing Gpu* fail-closed; prove
does not require a GPU; \texttt{measured\_j=None} unless Metered. & Soft
/ optional detect & Unmetered \\
P23c & wgpu / Metal tiny kernel & \texttt{claim.mol.kernel\_vector\_add}
& Clean-room WGSL vector-add; soft stub offline (same checksum); Gpu*
stamps \texttt{kernel:vector\_add} + \texttt{execution\_proof}; optional
overlapping meter may stamp package \texttt{measured\_j} with
\texttt{energy\_honesty=measured} (field name); no rail-sum package
joules. & Live Metal optional; fixture OK for prove & Unmetered
(default) / Metered only if meter returns \\
P24 & Ecosystem e2e certify (one receipt) &
\texttt{claim.mol.ecosystem\_e2e} & Lane → fabric (Cpu software
reference) → WASM → energy labels (estimated fuel; optional
meter\_sample) → host invoke only with \texttt{GrantReceipt} →
commit\textbar refuse on \textbf{one} receipt; refuse-without-grant →
\texttt{grant\_receipt\_required}. & none & Unmetered / Estimated \\
\end{longtable}
}

Outside the harness: live energy counters (RAPL, NVML) as the default
close, live GPU detection, on-device certificate hardware, trained model
weights, and silicon meters for Gap cells. None is claimed.

\subsection{II.6. Related work}\label{ii.6.-related-work}

\textbf{Cascades and metareasoning.} Viola and Jones (2001) ran a
cascade of boosted classifiers in which cheap stages reject most inputs.
Russell and Wefald (1991) priced each computation by its expected value
to the decision net of its cost and stopped when none had positive net
value. Horvitz (1987) set the same trade for bounded reasoning. Anytime
algorithms trade computation for quality with an interruptible stop
(Dean and Boddy, 1988; Zilberstein, 1996). Model cascading sends an
input to a larger model only when a smaller one is unsure (Varshney and
Baral, 2022), and FrugalGPT stops when a learned scorer accepts the
answer (Chen, Zaharia and Zou, 2023). Mixture of Limits differs in three
respects: the gears include non-neural mechanisms ahead of any model;
the stop is the predicate \(C(z)\), not a learned score; and cost is
joules on a receipt, not dollars or latency.

\textbf{Spend past the answer.} Reasoning models spend tokens after
their first correct answer without changing it (Chen et al., 2024), and
a survey collects methods that stop earlier (Sui et al., 2025).
Speculative decoding (Leviathan et al., 2023), Mixture-of-Depths (Raposo
et al., 2024), quantization (Frantar et al., 2022; Lin et al., 2024) and
paged KV caches (Kwon et al., 2023) lower the energy of a Model gear
once it runs. They do not decide whether it runs.

\textbf{Formula recovery.} AI Feynman (Udrescu and Tegmark, 2020), SINDy
(Brunton et al., 2016) and PySR (Cranmer, 2023) recover closed forms
from data; AlphaGeometry (Trinh et al., 2024) and AlphaProof (Hubert et
al., 2025) pair a proposer with a symbolic or Lean verifier. Each
produces Formula or certified Solver gears that close covered tasks at
lower energy than a generator.

\textbf{Routing among models.} These systems choose which model answers.
None tries a non-neural gear first or refuses with a typed reason.

{\def\LTcaptype{none} % do not increment counter
\begin{longtable}[]{@{}
  >{\raggedright\arraybackslash}p{(\linewidth - 12\tabcolsep) * \real{0.1429}}
  >{\raggedright\arraybackslash}p{(\linewidth - 12\tabcolsep) * \real{0.1429}}
  >{\raggedright\arraybackslash}p{(\linewidth - 12\tabcolsep) * \real{0.1429}}
  >{\raggedright\arraybackslash}p{(\linewidth - 12\tabcolsep) * \real{0.1429}}
  >{\raggedright\arraybackslash}p{(\linewidth - 12\tabcolsep) * \real{0.1429}}
  >{\raggedright\arraybackslash}p{(\linewidth - 12\tabcolsep) * \real{0.1429}}
  >{\raggedright\arraybackslash}p{(\linewidth - 12\tabcolsep) * \real{0.1429}}@{}}
\toprule\noalign{}
\begin{minipage}[b]{\linewidth}\raggedright
System
\end{minipage} & \begin{minipage}[b]{\linewidth}\raggedright
Title
\end{minipage} & \begin{minipage}[b]{\linewidth}\raggedright
ID
\end{minipage} & \begin{minipage}[b]{\linewidth}\raggedright
Class
\end{minipage} & \begin{minipage}[b]{\linewidth}\raggedright
What it routes or accelerates
\end{minipage} & \begin{minipage}[b]{\linewidth}\raggedright
Energy or cost as reported
\end{minipage} & \begin{minipage}[b]{\linewidth}\raggedright
Difference
\end{minipage} \\
\midrule\noalign{}
\endhead
\bottomrule\noalign{}
\endlastfoot
\textbf{RouteLLM} & \emph{RouteLLM: Learning to Route LLMs with
Preference Data} (arXiv); ICLR 2025 venue title uses \emph{from
Preference Data} &
arXiv:\href{https://arxiv.org/abs/2406.18665}{2406.18665} · ICLR 2025 &
Preference router (strong↔weak LLM) & Query → strong or weak LLM & Cost
/ quality (API \$); reported by the authors & Model↔model; no
Lookup/Formula floor; no typed refuse-to-commit \\
\textbf{PEARL} (routing) & \emph{PEARL: Performance and energy aware
routing for LLMs} &
DOI:\href{https://doi.org/10.1016/j.future.2025.108218}{10.1016/j.future.2025.108218}
· FGCS 176:108218 & Energy-aware multi-LLM router & Query → LLM under
energy cap (EMM predicts energy) & Authors report energy-aware routing
(GPU/infra) & Still picks a model; Mixture of Limits may refuse before
any model \\
\textbf{GreenServ} & \emph{GreenServ: Energy-Efficient Context-Aware
Dynamic Routing for Multi-Model LLM Inference} &
arXiv:\href{https://arxiv.org/abs/2601.17551}{2601.17551} · ICPE 2026 &
Contextual bandit multi-LLM router & Query features → LLM; accuracy vs
\textbf{measured GPU energy} (Zeus) & \textbf{Metered} (authors' GPU Wh
via Zeus) on their LLM pool & Model pool routing; Mixture of Limits
cascade is non-neural first \\
\textbf{VoltanaLLM} & \emph{VoltanaLLM: Energy-Efficient and SLO-Aware
Disaggregated LLM Serving via Adaptive Frequency Control and State-Space
Routing} & arXiv:\href{https://arxiv.org/abs/2509.04827}{2509.04827} &
P/D-disagg serving + frequency + state-space route & Prefill/decode
instances + GPU frequency under TTFT/ITL SLOs & \textbf{Metered}
(authors; pyNVML on A100/GH200) up to \textasciitilde36.3\% E2E GPU
energy vs max-freq baseline (field serving result) & Routes
\textbf{instances of the same generative stack}; Mixture of Limits may
never open decode \\
\textbf{HCSpec} & \emph{HCSpec: Two-Tier Horizontal Cascade Speculative
Decoding for High-Efficiency Large Language Model Inference} &
DOI:\href{https://doi.org/10.18653/v1/2026.acl-long.353}{10.18653/v1/2026.acl-long.353}
· ACL 2026 & Speculative decoding (draft↔target) & Position-specialized
draft cascade → target verify & Latency speedup (vs EAGLE-3 / AR) &
Speeds generation; does not refuse generation on floors \\
\textbf{CAS-Spec} & \emph{CAS-Spec: Cascade Adaptive Self-Speculative
Decoding for On-the-Fly Lossless Inference Acceleration of LLMs} &
arXiv:\href{https://arxiv.org/abs/2510.26843}{2510.26843} · NeurIPS 2025
& Self-speculative cascade (DSIA drafts) & Draft stages from target
(sparsity/quant) + DyTC & Latency speedup \textasciitilde1.1×-2.3× AR
(authors); lossless tokens & Model-internal draft hierarchy ≠ Mixture of
Limits floors \\
\textbf{GCD} & \emph{Grammar-Constrained Decoding for Structured NLP
Tasks without Finetuning} &
arXiv:\href{https://arxiv.org/abs/2305.13971}{2305.13971} ·
DOI:\href{https://doi.org/10.18653/v1/2023.emnlp-main.674}{10.18653/v1/2023.emnlp-main.674}
· EMNLP 2023 & Grammar-constrained LLM decoding & Mask logits to CFG &
Unmetered / quality metrics & Constrains \textbf{tokens}; Mixture of
Limits constrains \textbf{commit} \\
\end{longtable}
}

\subsection{II.7. Baselines and
falsifiers}\label{ii.7.-baselines-and-falsifiers}

A cascade claim is tested on one typed task set as a frontier of joules
per task against closure rate. The arms are: always the large model;
always the small model; a learned router (Ong et al., 2024); a scored
cascade (Chen, Zaharia and Zou, 2023); a selector that estimates
\(p_i / c_i\) per task; and an oracle that knows the cheapest sufficient
gear for each task. The fixed Lookup → Formula → Solver → Model order
claims the frontier only where its curve lies below every arm. The order
is falsified on a task class where measured \(p_i / c_i\) is not
decreasing along it, by the derivation of Section II.2. The harness
claims are falsified if \texttt{mol\ prove} does not exit 0 on the
stated revision.

\subsection{II.8. Limits}\label{ii.8.-limits}

The harness is constructive existence in software. No Part II figure is
a meter reading: catalog and Landauer energies are estimates, and
\texttt{measured\_j} is unset on the harness path. The optimal-order
result assumes independent gear outcomes; correlated failures change the
expected energy and require the full sum of Section II.2 with
conditional probabilities. Pass probabilities \(p_i\) must be estimated
per task class, and estimation error moves the order. The frontier
comparison of Section II.7 has not been run.

\subsection{II. References}\label{ii.-references}

Bennett, C. H. (1973). Logical reversibility of computation. \emph{IBM
Journal of Research and Development} 17(6):525-532.
https://doi.org/10.1147/rd.176.0525

Brunton, S. L., Proctor, J. L., and Kutz, J. N. (2016). Discovering
governing equations from data by sparse identification of nonlinear
dynamical systems. \emph{PNAS} 113(15):3932-3937.
https://doi.org/10.1073/pnas.1517384113

Chen, L., Zaharia, M., and Zou, J. (2023). FrugalGPT. arXiv:2305.05176.

Chen, X., et al.~(2024). Do NOT Think That Much for 2+3=? On the
Overthinking of o1-Like LLMs. arXiv:2412.21187.

CAS-Spec (2025). Cascade Adaptive Self-Speculative Decoding. NeurIPS.
arXiv:2510.26843.

Cranmer, M. (2023). Interpretable Machine Learning for Science with PySR
and SymbolicRegression.jl. arXiv:2305.01582.

Dean, T., and Boddy, M. (1988). An analysis of time-dependent planning.
AAAI-88, 49-54.

Fedus, W., Zoph, B., and Shazeer, N. (2022). Switch Transformers.
\emph{JMLR} 23(120).

Frantar, E., et al.~(2022). GPTQ. arXiv:2210.17323.

Geng, S., et al.~(2023). Grammar-Constrained Decoding for Structured NLP
Tasks without Finetuning. EMNLP.
https://doi.org/10.18653/v1/2023.emnlp-main.674

GreenServ (2026). Energy-Efficient Context-Aware Dynamic Routing for
Multi-Model LLM Inference. ICPE. arXiv:2601.17551.

HCSpec (2026). Two-Tier Horizontal Cascade Speculative Decoding. ACL.
https://doi.org/10.18653/v1/2026.acl-long.353

Horowitz, M. (2014). Computing's energy problem (and what we can do
about it). ISSCC. https://doi.org/10.1109/ISSCC.2014.6757323

Horvitz, E. J. (1987). Reasoning about beliefs and actions under
computational resource constraints. UAI. arXiv:1304.2759.

Howard, R. A. (1966). Information value theory. \emph{IEEE Transactions
on Systems Science and Cybernetics} 2(1):22-26.
https://doi.org/10.1109/TSSC.1966.300074

Hubert, T., et al.~(2025). Olympiad-level formal mathematical reasoning
with reinforcement learning. \emph{Nature}.
https://doi.org/10.1038/s41586-025-09833-y

Kwon, W., et al.~(2023). Efficient Memory Management for Large Language
Model Serving with PagedAttention. arXiv:2309.06180.

Landauer, R. (1961). Irreversibility and heat generation in the
computing process. \emph{IBM Journal of Research and Development}
5(3):183-191. https://doi.org/10.1147/rd.53.0183

Leviathan, Y., Kalman, M., and Matias, Y. (2023). Fast Inference from
Transformers via Speculative Decoding. ICML.

Lin, J., et al.~(2024). AWQ. MLSys. arXiv:2306.00978.

Ong, I., et al.~(2024). RouteLLM: Learning to Route LLMs with Preference
Data. arXiv:2406.18665.

PEARL (2025). Performance and energy aware routing for LLMs.
\emph{Future Generation Computer Systems} 176:108218.
https://doi.org/10.1016/j.future.2025.108218

Raposo, D., et al.~(2024). Mixture-of-Depths. arXiv:2404.02258.

Russell, S., and Wefald, E. (1991). Principles of metareasoning.
\emph{Artificial Intelligence} 49:361-395.
https://doi.org/10.1016/0004-3702(91)90015-C

Shazeer, N., et al.~(2017). Outrageously Large Neural Networks: The
Sparsely-Gated Mixture-of-Experts Layer. ICLR. arXiv:1701.06538.

Simon, H. A., and Kadane, J. B. (1975). Optimal problem-solving search:
all-or-none solutions. \emph{Artificial Intelligence} 6(3):235-247.
https://doi.org/10.1016/0004-3702(75)90002-8

Snell, C., Lee, J., Xu, K., and Kumar, A. (2024). Scaling LLM Test-Time
Compute Optimally can be More Effective than Scaling Model Parameters.
arXiv:2408.03314.

Sui, Y., et al.~(2025). Stop Overthinking: A Survey on Efficient
Reasoning for Large Language Models. arXiv:2503.16419.

Trinh, T. H., et al.~(2024). Solving olympiad geometry without human
demonstrations. \emph{Nature} 625:476-482.
https://doi.org/10.1038/s41586-023-06747-5

Udrescu, S.-M., and Tegmark, M. (2020). AI Feynman. \emph{Science
Advances} 6(16). https://doi.org/10.1126/sciadv.aay2631

Varshney, N., and Baral, C. (2022). Model cascading. EMNLP.
https://doi.org/10.18653/v1/2022.emnlp-main.756

Viola, P., and Jones, M. (2001). Rapid object detection using a boosted
cascade of simple features. CVPR.
https://doi.org/10.1109/CVPR.2001.990517

VoltanaLLM (2025). Energy-Efficient and SLO-Aware Disaggregated LLM
Serving. arXiv:2509.04827.

Zilberstein, S. (1996). Using anytime algorithms in intelligent systems.
\emph{AI Magazine} 17(3).

\subsection{II. Appendix.
Reproducibility}\label{ii.-appendix.-reproducibility}

\begin{Shaded}
\begin{Highlighting}[]
\ExtensionTok{cargo}\NormalTok{ test }\AttributeTok{{-}{-}workspace}
\ExtensionTok{cargo}\NormalTok{ run }\AttributeTok{{-}p}\NormalTok{ mol{-}cli }\AttributeTok{{-}{-}}\NormalTok{ prove   }\CommentTok{\# expect VERIFIED lines, exit 0}
\end{Highlighting}
\end{Shaded}


\end{document}
